\documentclass[12pt]{article}

\usepackage[a4paper, margin=1in]{geometry}
\usepackage{amsmath,amssymb,amsthm}
\usepackage{mathtools}
\usepackage{enumitem}
\usepackage{booktabs}
\usepackage[hidelinks]{hyperref}
\usepackage{fontspec}

\defaultfontfeatures{Ligatures=TeX}
\IfFontExistsTF{Times New Roman}
  {\setmainfont{Times New Roman}}
  {\setmainfont{TeX Gyre Termes}}
\IfFontExistsTF{Menlo}
  {\setmonofont{Menlo}}
  {\setmonofont{Latin Modern Mono}}

\newtheorem{theorem}{Theorem}[section]
\newtheorem{proposition}[theorem]{Proposition}
\newtheorem{lemma}[theorem]{Lemma}
\newtheorem{corollary}[theorem]{Corollary}
\newtheorem{conjecture}[theorem]{Conjecture}
\theoremstyle{definition}
\newtheorem{definition}[theorem]{Definition}
\newtheorem{example}[theorem]{Example}
\theoremstyle{remark}
\newtheorem{remark}[theorem]{Remark}

\newcommand{\Khat}{\widehat K}
\newcommand{\Glin}{G^{\mathrm{lin}}}
\newcommand{\dT}{\delta\langle T\rangle}
\newcommand{\ev}[1]{\langle #1\rangle}

\title{%
  \textbf{Higher-Dimensional Consistency of Geometry and Modular Flow:}\\
  \textbf{Physical Source Identification and a Second-Order Extension}%
}

\author{
  Sidong Liu, PhD \\
  \small iBioStratix Ltd \\
  \small \texttt{sidongliu@hotmail.com}
}

\date{October 2026}
\hypersetup{
  pdftitle={Higher-Dimensional Consistency of Geometry and Modular Flow: Physical Source Identification and a Second-Order Extension},
  pdfauthor={Sidong Liu}
}

\begin{document}
\emergencystretch=2em
\raggedbottom
\hbadness=5000
\vbadness=5000

\maketitle

\begin{abstract}
An observer reads ball algebras at different positions, sizes and rest
frames. We ask when those readings agree with one geometry. In a conformal
theory on Minkowski space, the first variation of the geodesic ball's
boundary area minus $(d-1)/R$ times its volume equals minus the
CHM-weighted linearized Einstein tensor at every radius. The conditional
law $\eta\,\delta F_B+\delta\langle K_B\rangle=0$ is therefore equivalent
to $G^{\mathrm{lin}}=(2\pi/\eta)\delta\langle T\rangle$.

A nonzero-frequency source identifies this intrinsic linear response up
to normalization even when internal length scales are allowed. Local
coherent and neutral scalar preparations provide smooth finite-energy
Radon witnesses; a bounded pulse inside a permitted shared hull supplies
a more general positive-energy test, with local source regularity still
required. A scalar's physical modular energy includes a boundary charge.
Its subtraction, or the matched Wald density, is essential: raw energy
cannot satisfy the pure-area bulk law on all small balls.

One frame fixes only the normal constraint. Boost comparisons determine
the sourced tensor, with small-rapidity stress noise amplified as
$\beta^{-2}$. Free Weyl curvature is invisible at linear order.
Stress-containing selected algebras give the same sourced tensor but
different relative entropies. Under the supplied radial readout model,
conditional multipole information falls as $e^{-4k\log k}$.

We extend the local conditional law to second order. A quadratic
extrinsic-curvature correction gives the full second-order Einstein
residual from all small boosted balls. A bent slice of flat spacetime
shows why the correction or a geodesic-slice prescription is necessary.
For the prepared conformal scalar, changed propagation, metric variation
and $\xi G\Phi^2$ supply explicit nonzero next-order source terms.
The common cone, area density and consistency mechanism remain inputs.
Curved stress readings are not assumed to be curved modular Hamiltonians;
the common quantum source and nonlinear modular response form a strong,
testable extension rather than an established consequence of the linear law.
\end{abstract}

\clearpage
\tableofcontents
\clearpage

\section{Introduction}\label{sec:intro}

Consciousness~I asks whether a physical description can be fixed before
any perspective is specified~\cite{PaperI}. In its chiral conformal net,
the geometry and modular data of a standard Cut determine each other.
Transverse directions change the question. A higher-dimensional observer
has a family of ball algebras, with different centres, radii and rest
frames. Can one geometry account for all their readings?

For conformal matter in Minkowski spacetime these regions carry the
bulk Casini--Huerta--Myers (CHM) stress charge~\cite{CasiniHuertaMyers2011}.
We first construct its geometric partner rather than postulating a
two-derivative conserved tensor. For every smooth perturbation the
area--volume variation of a geodesic ball equals the weighted integral
of the linearized Einstein tensor. The Gauss equation and intrinsic
gauge invariance explain why that tensor appears. Comparing it with the
CHM charge gives a conditional same-spacetime Einstein law, without a
holographic dictionary. Its coefficient is supplied, and its physical
consistency mechanism remains to be explained.

The next question is whether an experiment identifies the response law.
A homogeneous source does not. A source with nonzero-frequency weight
and an open set of boost directions identifies the entire compact
intrinsic response at one tested radius, up to normalization. We retain
the proof for tempered Radon Fourier measures, and provide actual scalar
preparations with finite positive energy. A neutral pulse belongs to a
proper stable hull; a shared non-scalar hull with translation permission
can also supply a bounded positive-energy pulse. Control permission,
smooth local source readings and source-apparatus costs are distinct inputs.

The physical one-sided scalar energy includes a surface charge. Keeping
it in the ledger yields a raw-energy obstruction and a matched Wald
coefficient; dropping it would falsely identify a geometric response.
The resulting tensor is common to stress-containing selected algebras.
Their relative entropies differ, not the emitted stress. A single rest
frame leaves momentum and spatial stresses underdetermined, and no
linear ball reading fixes free gravitational radiation.

Second order is not just a finer version of the same construction.
Slice extrinsic curvature, changed source propagation, quantum reference
stress and curved modular flow become relevant. We derive an
extrinsic-corrected leading small-ball law and prove its two-jet
equivalence to the second-order Einstein equation. A flat bent-slice
counterexample and an explicit next-order scalar source test it.
Section~\ref{sec:secondorder} states the remaining common-geometry
extension with its source, readout and intervention conditions.

The linear results require a conformal net, the bulk CHM datum, and a
common Minkowski cone. Q-RAIF~A studies conditions for common propagation
cones~\cite{QRA}; this paper does not derive the cone of its boosted family.
Its exact conformal net is a supplied fixed-point theory. Approximate
cones in a lattice realization do not by themselves establish the
exact CHM modular flow used here; nonconformal corrections require a
separate limit estimate \cite{CGM2016}.
The area density $\eta$ and modular consistency are also inputs.
Additional edge energy must be unchanged, separately read and subtracted,
or directly bypassed by bulk stress readings. At second order we supply
a corrected local \emph{stress} consistency condition, not a curved CHM
formula. Neither the Gauss identity nor stress conservation makes that
condition physically true by itself.

Sections~\ref{sec:modular}--\ref{sec:einstein} establish the linear law.
Sections~\ref{sec:access}--\ref{sec:selection} supply preparations,
identification and partial-observer limits; Section~\ref{sec:secondorder}
extends the geometry and source. Section~\ref{sec:multipoles} prices
radial resolution, and the conclusion returns to the series' question
of access and a common world. Indexed appendices retain the coefficient,
spectral-measure, source and entropy calculations.

\section{The modular family of balls}\label{sec:modular}

\subsection{Balls, diamonds and their modular Hamiltonians}\label{subsec:chm}

Work in Minkowski space \(\mathbb R^{1,d}\), \(d\ge2\), with
\(\eta=\mathrm{diag}(-1,1,\dots,1)\). A ball \(B=B(x_0,R,u)\) is fixed by a centre \(x_0\), a
radius \(R\) and a future unit timelike vector \(u\): it is the set of points \(x\) with
\(u\cdot(x-x_0)=0\) and \(|x-x_0|<R\), an open ball in the hyperplane \(\Sigma_u(x_0)\)
orthogonal to \(u\). Its causal diamond is \(D(B)\). We write
\[
w_R(y)=\frac{R^2-|y-x_0|^2}{2R},\qquad T_{uu}=T_{\mu\nu}u^\mu u^\nu ,
\]
and \(\Delta_u\) for the Laplacian of \(\Sigma_u\).

Let \(\mathcal A\) be a conformal net on \(\mathbb R^{1,d}\) with vacuum \(\Omega\). The
modular group of \((\mathcal A(D(B)),\Omega)\) acts geometrically, by the flow of the
conformal Killing vector that preserves \(D(B)\)~\cite{HislopLongo1982,BGL1993}. In
coordinates with \(x_0=0\) and \(u=\partial_t\),
\begin{equation}\label{eq:zeta}
\zeta_B=\frac1{2R}\Big[(R^2-t^2-r^2)\,\partial_t-2t\,x^i\partial_i\Big],
\end{equation}
normalized so that its surface gravity at the edge \(\partial B\) is one; the modular period
is then \(2\pi\), as for the boost of a wedge. The one-sided CHM modular energy is the stress charge
\begin{equation}\label{eq:chm}
K_B=2\pi\int_\Sigma T_{\mu\nu}\zeta_B^\mu n^\nu\,dV
=2\pi\int_B w_R T_{uu}\,d^dx .
\end{equation}
Here \(\Sigma\) is a Cauchy surface of the diamond and the second expression uses its
central ball. In a theory without an additional edge charge, a faithful type~I
regulator identifies it with \(H_B=-\log\rho_B\), up to the additive normalization
constant \cite{CasiniHuertaMyers2011}. For a free conformal scalar the physical
one-sided energy includes a surface term \cite{CasiniMazzitelliTeste2015};
Section~\ref{subsec:prepared} keeps that term separate. In the continuum \(K_B\)
denotes the bulk charge and its
vacuum-subtracted expectations, with the required stress-smearing limits assumed finite.

The full modular generator is a different object:
\(\Khat_B=-\log\Delta_B\). It annihilates the vacuum and implements modular covariance.
At a regulated bipartition,
\[
\Khat_B=H_B\otimes1-1\otimes H_{B^c}^{\mathsf T}.
\]
Thus it agrees with the physical one-sided energy in expectation differences for a
unitary local to \(D(B)\), whose complement state is unchanged. It agrees with the
bulk charge as well when any boundary contribution is unchanged. Arbitrary state variations need not
have this property. Throughout, \(\delta\ev{K_B}\) means the vacuum-subtracted CHM charge,
and modular covariance uses \(\Khat_B\).

\subsection{Two kinds of variation}\label{subsec:variations}

HAFF-D separates two ways in which a modular Hamiltonian can change~\cite{HAFFD}, and both
occur in the ball family.

A \emph{state variation} keeps the region and changes the state:
\(\delta\ev{K_B}=2\pi\int_Bw_R\,\dT_{uu}\). For a unitary \(W\in\mathcal A(D(B))\)
with unchanged boundary charge, Araki's
relative entropy of \(\omega_W=\ev{W\Omega,\cdot\,W\Omega}\) is
\begin{equation}\label{eq:relent-unitary}
S_B(\omega_W\Vert\omega)=\ev{W\Omega,\Khat_BW\Omega}=\delta\ev{K_B}_{\omega_W},
\end{equation}
which Consciousness~I uses in \(1+1\) dimensions~\cite{PaperI,Araki1976}. For a first-order
perturbation \(\rho=\sigma+\epsilon\,\delta\rho\) of a regulated (type~I) version, the
first law is \(\delta S=\delta\ev{H_B}\); it equals the bulk
\(\delta\ev{K_B}\) when the boundary variation is zero.

A \emph{region variation} keeps the state and moves the ball. The family of all diamonds
is one orbit of the conformal group, of dimension \(2(d+1)\): \(d+1\) for the centre, one
for the radius and \(d\) for the rest frame. A one-parameter conformal group
\(e^{-iaP_\xi}\) moves \(D(B)\) to \(D(B_a)\) and
\(\Khat_{B_a}=e^{-iaP_\xi}\Khat_Be^{iaP_\xi}\), so
\begin{equation}\label{eq:regionmove}
\delta\Khat_B=-i\,\delta a\,[P_\xi,\Khat_B],\qquad \ev{\Omega,\delta\Khat_B\,\Omega}=0 ,
\end{equation}
because \(P_\xi\Omega=0\) and \(\Khat_B\Omega=0\). This is the variation formula of HAFF-D for a
half-sided modular family. Null translations of a wedge give half-sided modular
inclusions~\cite{Wiesbrock1993,Borchers1992}, and balls are conformal images of wedges, so
the supplied ball family has the region interface required by HAFF-D's conditional holographic matching. Vacuum covariance concerns the full generator.

This transitivity concerns a single diamond, not an arbitrary pair.
A nested pair of balls carries the inversive-distance modulus; it is not
determined by a dimension count for the orbit of one ball.

\subsection{Conservation from the diamond}\label{subsec:conservation}

The algebra of the diamond is the same for every Cauchy surface of \(D(B)\).
For its bulk stress charge, independence of the first form of~\eqref{eq:chm}
from \(\Sigma\) is exactly conservation and
tracelessness.

\begin{proposition}[Conservation from the diamond]\label{prop:conservation}
Let \(T_{\mu\nu}\) be a smooth symmetric tensor field on an open set
\(U\subset\mathbb R^{1,d}\). The charge \(Q_\Sigma(B)=\int_\Sigma T_{\mu\nu}\zeta_B^\mu n^\nu
dV\) is the same for all Cauchy surfaces \(\Sigma\) of \(D(B)\), for every diamond with
\(\overline{D(B)}\subset U\), if and only if \(\partial^\mu T_{\mu\nu}=0\) and
\(T^\mu{}_\mu=0\) in \(U\).
\end{proposition}

\begin{proof}
\(\zeta_B\) is a conformal Killing vector,
\(\partial_\mu\zeta_\nu+\partial_\nu\zeta_\mu=\frac{2}{d+1}(\partial\cdot\zeta_B)\eta_{\mu\nu}\),
with \(\partial\cdot\zeta_B=-(d+1)t/R\) in the coordinates of~\eqref{eq:zeta}. Hence
\[
\partial^\mu\big(T_{\mu\nu}\zeta_B^\nu\big)
=\zeta_B^\nu\,\partial^\mu T_{\mu\nu}+\frac{\partial\cdot\zeta_B}{d+1}\,T^\mu{}_\mu .
\]
Two Cauchy surfaces of \(D(B)\) share the edge \(\partial B\) and bound a region of
\(D(B)\); if the right side vanishes, Stokes' theorem gives equal charges. Conversely,
deforming \(\Sigma\) near an interior point \(x\) shows that the left side vanishes at
every \(x\in D(B)\), for every diamond. As \(B\) varies over diamonds containing a fixed
\(x\), the pairs \((\zeta_B(x),\partial\cdot\zeta_B(x))\) span
\(\mathbb R^{d+1}\times\mathbb R\), so both terms vanish separately.
\end{proof}

For the anomaly-free flat conformal stress both properties hold.
The proposition identifies what slice-independent stress-charge representation requires: a modular datum is attached to a region's algebra, not
to a slice; provided this same bulk stress charge represents it on every
Cauchy surface, the stated equality forces conservation and a vanishing trace. Conservation is therefore not an observation to be placed beside the Bianchi
identity; it is a property of the modular family.

The finite-radius charge has a controlled expansion with leading
coefficient $c_0=V_d/(d+2)$; its exact Fourier transfer and all radial
coefficients are given in Appendix~\ref{subsec:series}. In particular,
an open set of normals recovers a smooth source tensor from its leading
small-ball charges (Corollary~\ref{cor:recovery}).

\section{The geometric functional of a ball}\label{sec:functional}

\subsection{What a ball functional may read}\label{subsec:balldef}

Let \(h_{\mu\nu}\) be a smooth metric perturbation of Minkowski space, and \(h_{ij}\) the
perturbation of the metric induced on \(\Sigma_u(x_0)\).

\begin{definition}[Ball functional]\label{def:ballfunctional}
A \emph{ball functional} assigns to every ball \(B=B(x_0,R,u)\) and every \(h\) a number
\(\delta F_B[h]\) such that:
\begin{enumerate}[label=(F\arabic*),leftmargin=3em,itemsep=1pt]
\item it is linear in \(h\);
\item it depends only on the induced perturbation \(h_{ij}\) on \(\overline B\);
\item it is unchanged by \(h_{ij}\mapsto h_{ij}+\partial_i\xi_j+\partial_j\xi_i\) for every
smooth vector field \(\xi\) on a neighbourhood of \(\overline B\) in \(\Sigma_u\);
\item it is invariant under the orthogonal group of \(\Sigma_u\) fixing \(x_0\);
\item it has the same form for every centre and every rest frame, its dependence on \(R\)
and on any constants being arbitrary;
\item it is a distribution of finite order supported in \(\overline B\).
\end{enumerate}
\end{definition}

No scale covariance is assumed, so the functional may contain any number of length scales
besides \(R\). The restrictive condition is (F2): the functional reads the intrinsic
geometry of the ball. Three considerations support it. The region to which the modular
datum belongs is the ball on its slice, whose geometric data are the induced metric and
the extrinsic curvature. The extrinsic curvature reverses sign under reflection of time
about \(\Sigma_u\), while \(\delta\ev{K_B}\) does not change when the state is reflected in
a theory with that symmetry; a linear functional matched to the modular energy for all
states cannot depend on it. And second time derivatives of \(h\) are not data of the region
at all: reading them presupposes an evolution law for the metric, which is what one wants
to obtain. Proposition~\ref{prop:failures} shows what happens when (F2) is dropped.

\begin{lemma}[Ball functionals see only the Einstein tensor]\label{lem:tau}
Let \(\delta F\) satisfy (F1)--(F4) and (F6), and let
\(\tau_u=\partial_i\partial_jh_{ij}-\Delta_uh_{kk}\) be the linearized scalar curvature of
\(\Sigma_u\).
\begin{enumerate}[label=(\alph*),leftmargin=2em,itemsep=1pt]
\item For a plane wave \(h_{ij}=\mathrm{Re}(\varepsilon_{ij}e^{ip\cdot x})\) on \(\Sigma_u\),
\(\delta F_B[h]=\mathrm{Re}\big(m_B(|p|)\,\hat\tau\,e^{ip\cdot x_0}\big)\) with
\(\hat\tau=p^2\varepsilon_{kk}-p_ip_j\varepsilon_{ij}\) and \(m_B\) an even entire function,
the \emph{transfer function} of the functional.
\item For polynomial \(h\), \(\delta F_B[h]=\sum_k m_k\,\Delta_u^k\tau_u(x_0)\), where
\(m_B(q)=\sum_km_k(-q^2)^k\).
\item \(\tau_u=2\,\Glin_{uu}\).
\item There is a rotation-invariant scalar distribution \(\mu_B\), supported in
\(\overline B\), such that for every smooth \(h\) on a neighbourhood of that ball,
\[
\delta F_B[h]=\langle\mu_B,\tau_u[h]\rangle,
\qquad \widehat\mu_B(p)=m_B(|p|).
\]
No growth condition on \(h\) outside the ball is required.
\end{enumerate}
\end{lemma}

The proof is in Appendix~\ref{app:lem:tau}.


At linear order about a flat slice, the Gauss equation equates half
the intrinsic scalar curvature with the normal--normal Einstein component.
The quadratic extrinsic term is essential beyond that order
(Section~\ref{sec:secondorder}). The
Bianchi identity, which makes \(\Glin\) divergence-free, therefore comes with the
functional; it is not a separate premise.

\begin{corollary}[The Weyl curvature is invisible]\label{cor:weyl}
If \(h\) is smooth and \(\Glin=0\), every ball functional vanishes on \(h\) for every
ball. In particular, for \(d\ge3\) a free gravitational wave, whose Riemann tensor is pure
Weyl curvature, gives zero for every ball, every rest frame and every functional satisfying
(F1)--(F4) and (F6).
\end{corollary}

\subsection{The functional}\label{subsec:identity}

For a ball \(B=B(x_0,R,u)\) and a metric \(g=\eta+\epsilon h\), let \(B_\epsilon\) be the
geodesic ball of radius \(R\) about \(x_0\) in the metric induced on \(\Sigma_u(x_0)\),
\(A_\epsilon\) the area of its boundary and \(V_\epsilon\) its volume. Define
\begin{equation}\label{eq:F}
F_B[g]=A(\partial B_\epsilon)-\frac{d-1}{R}\,V(B_\epsilon),\qquad
\delta F_B[h]=\frac{d}{d\epsilon}F_B[\eta+\epsilon h]\Big|_{\epsilon=0}.
\end{equation}
For \(d=2\) the area is the length of the boundary circle.

\begin{theorem}[Exact identity]\label{thm:identity}
For every smooth \(h\) and every ball,
\begin{equation}\label{eq:identity}
\delta F_B[h]=-\frac12\int_B w_R\,\tau_u\,d^dx=-\int_B w_R\,\Glin_{uu}\,d^dx .
\end{equation}
\end{theorem}

\begin{proof}
Put \(x_0=0\) and write points of \(\Sigma_u\) as \(s\hat n\). To first order the geodesic
distance from \(0\) to \(s\hat n\) is \(s+\frac12\int_0^sh_{nn}(\sigma\hat n)\,d\sigma\),
because the length functional is stationary on the straight line. The boundary of
\(B_\epsilon\) is therefore \(s=R-\epsilon\Phi(\hat n)\) with
\(\Phi(\hat n)=\frac12\int_0^Rh_{nn}(\sigma\hat n)\,d\sigma\). The area element of a surface
changes by half the trace of \(h\) over its tangent space, and the volume element by half
the full trace, so
\[
\delta A=\frac12\oint_{S_R}\big(h_{kk}-h_{nn}\big)\,dA-(d-1)R^{d-2}\oint\Phi\,d\Omega,\qquad
\delta V=\frac12\int_{B}h_{kk}\,d^dx-R^{d-1}\oint\Phi\,d\Omega .
\]
In \(\delta A-\frac{d-1}{R}\delta V\) the terms in \(\Phi\) cancel:
\[
\delta F_B=\frac12\oint_{S_R}\big(h_{kk}-h_{nn}\big)\,dA-\frac{d-1}{2R}\int_Bh_{kk}\,d^dx .
\]
On the other side use \(w_R=0\) and \(\partial_nw_R=-1\) on \(S_R\),
\(\partial_iw_R=-x_i/R\) and \(\partial_i\partial_jw_R=-\delta_{ij}/R\). Integrating by parts
twice,
\(\int_Bw_R\,\partial_i\partial_jh_{ij}=\oint h_{nn}\,dA-\frac1R\int_Bh_{kk}\) and
\(\int_Bw_R\,\Delta h_{kk}=\oint h_{kk}\,dA-\frac dR\int_Bh_{kk}\). Subtracting gives
\(\int_Bw_R\tau_u=\oint(h_{nn}-h_{kk})\,dA+\frac{d-1}R\int_Bh_{kk}\), which is
\(-2\,\delta F_B\). The second equality in~\eqref{eq:identity} is
Lemma~\ref{lem:tau}(c).
\end{proof}

The geometric identity can also be evaluated without the linearized formulas, by
shooting geodesics with their
Jacobi fields in the metric \(\delta+\epsilon h\) and differencing in \(\epsilon\).
For a non-conformally-flat perturbation, this evaluation agrees with
the right side of~\eqref{eq:identity} to
relative accuracy \(2\times10^{-11}\) in \(d=2\) and \(10^{-10}\) in \(d=3\).

The content of Theorem~\ref{thm:identity} is known: it is the flat-space, linear
geometric part of the first law of causal diamonds of Jacobson and
Visser~\cite{JacobsonVisser2019}, there obtained from the covariant phase space of general
relativity with the Einstein constraint imposed. The proof above uses no field equation.
That is what the next section needs: \eqref{eq:identity} holds for every perturbation,
whatever its dynamics.

\subsection{Where the coefficient comes from}\label{subsec:lambda}

The volume term in~\eqref{eq:F} is the fixed-volume prescription of Jacobson's
entanglement-equilibrium argument~\cite{Jacobson2016}, where it is motivated physically.
In the modular family it has a structural origin.

\begin{proposition}[The coefficient]\label{prop:lambda}
Let \(F_\lambda=A-\lambda V\) for geodesic balls.
\begin{enumerate}[label=(\alph*),leftmargin=2em,itemsep=1pt]
\item Fix a reference radius \(R_0\) and hold \(\lambda_0\) fixed when varying \(R\).
Then
\[
\left.\frac{d}{dR}\left[\Omega_{d-1}
 \left(R^{d-1}-\frac{\lambda_0R^d}{d}\right)\right]\right|_{R=R_0}=0
\quad\Longleftrightarrow\quad \lambda_0=\frac{d-1}{R_0}.
\]
Centre and frame moves leave the flat value unchanged for every \(\lambda_0\).
This fixes the Lagrange multiplier for the variation at \(R_0\).
\item Write \(\delta F_\lambda=\sum_k\gamma_k(\lambda)\,R^{d+1+2k}\Delta_u^k\tau_u(x_0)\) for
polynomial \(h\). Then
\[
\frac{\gamma_k(\lambda)}{-c_k/2}=\frac{d+2k+2-\lambda R}{2k+3},
\]
which is independent of \(k\) if and only if \(\lambda=(d-1)/R\). Thus \(\delta F_\lambda\)
has the modular weight \(w_R\) at every order in \(R\) exactly for this coefficient.
\end{enumerate}
\end{proposition}

The proof is in Appendix~\ref{app:prop:lambda}.


Agreement of a geometric functional with the modular data at leading order in \(R\) does
not fix \(\lambda\); only the coefficient of the Einstein equation would change. Matching the full kernel fixes it to the fixed-reference multiplier.
The family \(F_R=A_R-(d-1)V_R/R\) itself is not stationary:
\(F_R[\eta]=\Omega_{d-1}R^{d-1}/d\), so its radius derivative is
\(\Omega_{d-1}(d-1)R^{d-2}/d\). These two variations must be distinguished.

\section{Modular consistency and the linearized Einstein equation}\label{sec:einstein}

\begin{definition}[Modular consistency]\label{def:MC}
Let a state have smooth stress-tensor expectation \(\dT\), let \(h\) be a metric
perturbation, \(\delta F\) a ball functional and \(\eta>0\) a constant. The pair satisfies
\((\mathrm{MC}_\eta)\) on a family of balls if
\begin{equation}\label{eq:MC}
\eta\,\delta F_B[h]+\delta\ev{K_B}=0
\end{equation}
for every ball of the family.
\end{definition}

Here \(K_B\) is the bulk datum defined by~\eqref{eq:chm}. If the physical one-sided
modular energy is \(H_B=K_B+Q_B\), the same condition uses
\(\delta\ev{H_B}-\delta\ev{Q_B}\). It uses the raw physical energy only when the
boundary variation vanishes.

The constant \(\eta\) has the dimension of an inverse area. For first-order perturbations
with unchanged boundary charge, where \(\delta S_B=\delta\ev{K_B}\), and with
\(\delta F=\delta F_B\) of~\eqref{eq:F},
condition~\eqref{eq:MC} says that \(\eta A+S\) is stationary at fixed volume, which is
Jacobson's entanglement equilibrium~\cite{Jacobson2016}. We keep the modular energy, not
the entropy, in~\eqref{eq:MC}; Remark~\ref{rem:energy} and Section~\ref{sec:selection}
explain why. In words, \((\mathrm{MC}_\eta)\) says that the modular energy of every ball is
balanced by one geometric response of the ball, with one constant.
If a boundary charge varies, the first-order entropic version uses
\(\delta S_B-\delta\ev{Q_B}\), not the uncorrected entropy variation.

\subsection{Equivalence}\label{subsec:equivalence}

\begin{theorem}[Equivalence]\label{thm:equivalence}
Let \(U\subset\mathbb R^{1,d}\) be open, and let \(\dT\) and \(h\) be smooth on \(U\). With
\(\delta F\) the functional~\eqref{eq:F}, the following are equivalent.
\begin{enumerate}[label=(\alph*),leftmargin=2em,itemsep=1pt]
\item \((\mathrm{MC}_\eta)\) holds for all balls \(\overline B\subset U\), in all rest frames and at all
radii.
\item \((\mathrm{MC}_\eta)\) holds to leading order in \(R\) at every \(x_0\in U\), for all
\(u\) in some nonempty open set of rest frames.
\item \(\Glin_{\mu\nu}=\dfrac{2\pi}{\eta}\,\dT_{\mu\nu}\) in \(U\).
\end{enumerate}
\end{theorem}

\begin{proof}
(c)\(\Rightarrow\)(a): by Theorem~\ref{thm:identity},
\(\eta\,\delta F_B=-\eta\int_Bw_R\Glin_{uu}=-2\pi\int_Bw_R\dT_{uu}=-\delta\ev{K_B}\). (a)
\(\Rightarrow\)(b) is immediate. (b)\(\Rightarrow\)(c): the leading terms of both sides
of~\eqref{eq:MC} carry the same factor \(V_dR^{d+1}/(d+2)\)
(Proposition~\ref{prop:series} and Theorem~\ref{thm:identity}), so
\(\eta\Glin_{uu}(x_0)=2\pi\dT_{uu}(x_0)\) for an open set of \(u\), and
Corollary~\ref{cor:recovery} gives the tensor equation.
\end{proof}

\emph{Normalization.} The factor \(V_d/(d+2)\) appears on both sides and cancels, so the
coupling is \(2\pi/\eta\) in every dimension. With \(\eta=1/4G\) it is \(8\pi G\): the
\(2\pi\) is the modular period of a ball and the \(1/4\) is the Bekenstein--Hawking area
density.

\emph{All radii from leading order.} Agreement at leading order already forces the field
equation, and the field equation makes the two sides agree at every radius. The finite-size
corrections of the modular side and of the geometric side are therefore identical, term by
term in~\eqref{eq:series}; for a Fourier mode both have the transfer function \(f\)
of~\eqref{eq:transfer}. Given the field equation, the full ball-dependent
functional agrees with the modular data on every ball; this is the converse direction.

\emph{Conservation and trace.} The Einstein tensor is divergence-free by the Bianchi
identity and \(\dT\) is conserved and traceless by Proposition~\ref{prop:conservation}; the
equation is consistent with both and requires \(R^{\mathrm{lin}}=0\) in \(U\). Neither
property is an additional premise.

\emph{Background.} The background is Minkowski space, so no cosmological term appears.
For diamonds in de~Sitter or anti-de~Sitter space the first law of Jacobson and Visser
contains the cosmological constant and its variation~\cite{JacobsonVisser2019}; the
corresponding version of Theorem~\ref{thm:identity} is not derived here.

\subsection{Uniqueness}\label{subsec:uniqueness}

The canonical identity holds for every metric perturbation. Identifying that response
from data requires a source that excites more than the zero-frequency channel.

\begin{lemma}[Conserved amplitudes under frame variation]\label{lem:quadric}
Let \(k\ne0\), and let \(t,g\) be symmetric complex tensors with
\(k^\mu t_{\mu\nu}=k^\mu g_{\mu\nu}=0\), \(t\ne0\).
On a nonempty open set of future unit timelike vectors suppose
\(g_{uu}=a(q_u)t_{uu}\), where \(q_u^2=k^2+(u\cdot k)^2\) and \(a\) is analytic
on an interval of attained momenta. Then \(g=a_0t\) for a constant \(a_0\).
Tracelessness is not needed.
\end{lemma}

\begin{proof}
For timelike \(k\), choose its rest frame. Conservation sets \(g_{0a}=t_{0a}=0\).
Write \(u=(\sqrt{1+r^2},rn)\). At fixed \(r>0\), equality on an open patch of
the sphere implies \(g_{ij}=a(q_r)t_{ij}\) by equality of quadratic forms.
The nonzero tensor \(t\) fixes the same scalar as \(r\) varies.

For spacelike \(k\), put it along the first spatial coordinate. Conservation sets
\(g_{1a}=t_{1a}=0\). At fixed \(u^1=b\), the remaining coordinates run over the
timelike hyperboloid of squared norm \(-1-b^2\).
A homogeneous quadratic form vanishing on an open patch of this hyperboloid vanishes
on the open timelike cone by analytic continuation and homogeneity, hence identically.
Thus \(g=a(q_b)t\), with one scalar independent of \(b\).

For null \(k\), take \(k\propto(1,1,0,\ldots)\). Conservation gives
\(g_{0a}+g_{1a}=t_{0a}+t_{1a}=0\). At fixed \(r=u^0-u^1>0\), the transverse
coordinates \(v\) vary freely, \(q_u\) is fixed, and
\[
g_{uu}=g_{00}r^2+2r g_{0A}v^A+g_{AB}v^Av^B.
\]
Equality on an open patch is equality of these polynomials. It gives \(g=a(q_r)t\);
again \(t\ne0\) fixes the scalar. An open frame set contains patches with an interval
of \(r\) or \(b\), away from any isolated critical value of \(q_u\).
\end{proof}

\begin{theorem}[Nonzero-frequency identification]\label{thm:uniqueness}
Fix \(R>0\) and let \(\delta F\) obey (F1)--(F6), with transfer function \(m_R\).
Suppose
\[
\dT_{ab}(x)=\operatorname{Re}(t_{ab}e^{ik\cdot x}),\qquad
\Glin_{ab}[h](x)=\operatorname{Re}(g_{ab}e^{ik\cdot x}),
\]
where \(k\ne0\), \(t\ne0\), and \(t,g\) are conserved symmetric amplitudes.
If \((\mathrm{MC}_\eta)\) holds on a nonempty open set of centres and on a nonempty open set of rest
frames at this radius, then
\[
m_R(q)=-\frac c2\widehat w_R(q)\quad\hbox{for every }q,\qquad
g_{ab}=\frac{2\pi}{c\eta}t_{ab},\qquad c\in\mathbb R\setminus\{0\}.
\]
Consequently \(\delta F=c\,\delta F_B\) at this radius on arbitrary smooth metric inputs.
If the same witness satisfies consistency at further radii, the same \(c\) applies
at each tested radius. A hypothesis restricted to \(R<R_0\) identifies only that range.

More generally let the spacetime Fourier transforms of \(\dT\) and \(\Glin[h]\) be
tempered tensor-valued complex Radon measures, meaning that their total variations
are locally finite and have at most polynomial growth, with the source measure nonzero away
from \(k=0\). With consistency at all spacetime centres and an open frame set,
the response is identified at
this radius and \(\Glin[h]=(2\pi/(c\eta))\dT\) as distributions.
This includes continuous Fourier densities and sources stationary in one frame with
a continuous spatial spectrum. General tempered distributions need a separate argument.
\end{theorem}

The proof is in Appendix~\ref{app:thm:uniqueness}.


One radius is enough because changing the rest frame turns a nonzero spacetime
frequency into an interval of spatial frequencies, and the compact support of the
response supplies analytic continuation. Exact identification does not imply stable
identification from finitely many noisy frames. The theorem concerns response laws tested
by such data, rather than a physical derivation of the modular-consistency hypothesis.

\begin{corollary}[The spectral criterion is necessary and sufficient]\label{cor:spectraliff}
Let \(d\ge2\), fix a radius, and suppose a nonzero smooth source and a smooth metric
satisfy \((\mathrm{MC}_\eta)\) for some functional obeying (F1)--(F6), at every centre
and on an open set of frames. Suppose both tensor Fourier transforms belong to the
tempered Radon class of Theorem~\ref{thm:uniqueness}. The data identify the canonical
functional up to normalization if and only if the source measure has nonzero weight
away from zero. If it is supported at zero, both tensors are constant and the data
test only \(m_R(0)\).
\end{corollary}

The proof is in Appendix~\ref{app:cor:spectraliff}.


The nonzero-source hypothesis excludes a degenerate experiment: if both source and
Einstein tensor vanish, every admissible functional fits. Polynomial sources fall
outside the Radon class; their Fourier transforms contain derivatives of \(\delta_0\).
A polynomial tensor of degree \(n\) probes only the finitely many coefficients
\(m_k\) with \(2k\le n\), so it also leaves higher spatial derivatives undetermined.

\begin{example}[A homogeneous source leaves the response open]\label{ex:homogeneous}
In \(D=d+1\ge3\) spacetime dimensions put
\(T_{ab}=\operatorname{diag}(d,1,\ldots,1)\), \(\kappa=2\pi/\eta\),
\(x^2=\eta_{ab}x^ax^b\), and \(Txx=T_{ab}x^ax^b\). The quadratic metric
\[
h_{ab}=-\frac{\kappa}{2(D-2)}\bigl(T_{ab}x^2+\eta_{ab}Txx\bigr)
\]
has \(R^{\mathrm{lin}}=0\) and \(\Glin_{ab}=\kappa T_{ab}\).
For every \(\ell\ne0\), the intrinsic functional
\[
\delta F^{\mathrm{new}}_B[h]=\delta F_B[h]
+\ell^2\int_Bw_R\Delta_u\tau_u[h]\,d^dx
\]
obeys (F1)--(F6) and satisfies exactly the same consistency condition on this
pair at every centre, radius and frame. Its extra term vanishes because
\(\Glin\) is constant. It is not a multiple of the canonical functional:
its transfer function is \((-\tfrac12-\ell^2q^2)\widehat w_R(q)\).
A homogeneous source tests only the zero-momentum response.
\end{example}

The Einstein tensor appears through intrinsic slice curvature and the Gauss equation.
Nonzero-frequency data then identify its two-derivative response: no independent
assumption excluding internal length scales is needed. Actual source preparation is
therefore part of what makes a geometric law identifiable.

\subsection{What fails without each hypothesis}\label{subsec:failures}

\begin{proposition}\label{prop:failures}
\begin{enumerate}[label=(\alph*),leftmargin=2em,itemsep=2pt]
\item \emph{One order in \(R\) is not enough.} The area of the geodesic sphere at fixed
radius, \(\lambda=0\) in Proposition~\ref{prop:lambda}, satisfies \((\mathrm{MC}_\eta)\) to
leading order with \(\Glin=\frac{6\pi}{(d+2)\eta}\dT\), but its coefficient of order
\(R^{d+3}\) differs from the modular one by the factor \(1-\frac{2(d-1)}{5(d+2)}\)
(\(-16\%\) for \(d=3\)). Consistency at that order would require \(\Delta_u\dT_{uu}=0\) in
every frame. For a Fourier mode the two transfer functions differ by \(0.29\%\), \(1.2\%\),
\(5.2\%\) and \(14\%\) at \(R|k_u|=0.5,1,2,3\) (\(d=3\)).
\item \emph{Reading off the slice.} The functional
\(\delta F^{(\ell)}_B=-\int_Bw_R\,(\Glin_{uu}-\ell^2\Box\Glin_{uu})\) satisfies
(F1) and (F3)--(F6) but not (F2). \((\mathrm{MC}_\eta)\) holds for it at all radii and in all
frames if and only if \((1-\ell^2\Box)\Glin_{\mu\nu}=(2\pi/\eta)\dT_{\mu\nu}\), a
fourth-order equation.
\item \emph{Reading only the trace-free part.} The functional
\(-\int_Bw_R\,(R^{\mathrm{lin}}_{uu}+R^{\mathrm{lin}}/(d+1))\) gives
\(R^{\mathrm{lin}}_{\mu\nu}-\frac{1}{d+1}\eta_{\mu\nu}R^{\mathrm{lin}}=(2\pi/\eta)\dT_{\mu\nu}\);
the Bianchi identity makes \(R^{\mathrm{lin}}\) constant, and the result is
\(\Glin_{\mu\nu}+\Lambda\eta_{\mu\nu}=(2\pi/\eta)\dT_{\mu\nu}\) with an integration constant
\(\Lambda=\frac{d-1}{2(d+1)}R^{\mathrm{lin}}\).
\end{enumerate}
\end{proposition}

\begin{proof}
(a) From Proposition~\ref{prop:lambda}(b) with \(\lambda=0\) the ratios to the modular
coefficients are \((d+2)/3\) and \((d+4)/5\) at the first two orders; the transfer
functions are computed from the linearized geodesic ball. (b) and (c) follow as
in Theorem~\ref{thm:equivalence}, with the operator inside the integral in place of
\(\Glin_{uu}\); in (c), \(\partial^\mu R^{\mathrm{lin}}_{\mu\nu}=\frac12\partial_\nu
R^{\mathrm{lin}}\) and \(\partial^\mu\dT_{\mu\nu}=0\) give
\(\big(\frac12-\frac1{d+1}\big)\partial_\nu R^{\mathrm{lin}}=0\).
\end{proof}

Part~(b) is the general lesson. Higher-derivative laws arise when the geometric response of
a region depends on how the geometry evolves away from the region's slice; the analysis of
higher-curvature entanglement equilibrium by Bueno et al.\ fits this
pattern~\cite{BMSV2017}. Part~(c) shows where a cosmological constant could enter as an
integration constant: only through a functional blind to the trace, which the intrinsic one
is not.

\begin{remark}[Energy, not entropy]\label{rem:energy}
For a unitary \(W\in\mathcal A(D(B))\) with unchanged boundary charge,
the entropy of the ball does not change at a
regulator preserving that local unitary, while
\(\delta\ev{K_B}=S_B(\omega_W\Vert\omega)\ge0\),
with strict inequality whenever the local state changes,
by~\eqref{eq:relent-unitary}. For a coherent excitation strictly local to \(D(B)\), the
modular energy and the relative entropy of that ball are equal and of second order in
the amplitude. A condition written with entropy,
\(\eta\,\delta F_B+\delta S_B=0\), would then require \(\delta F_B=0\) on that ball,
in conflict with the semiclassical Einstein response to a nontrivial preparation.
Condition~\eqref{eq:MC} has no such problem. Entanglement equilibrium agrees with it for
perturbations whose modular energy is of first order, and only for those.
\end{remark}

\section{What an observer's ball data fix}\label{sec:access}

An observer does not have all balls. It has balls near its worldline, of sizes between a
resolution and a scale of interest, in the rest frames it can compare. This section asks
what such data fix through \((\mathrm{MC}_\eta)\), and with what stability.

\subsection{One rest frame}\label{subsec:singleframe}

\begin{proposition}[One rest frame gives the Hamiltonian constraint]\label{prop:singleframe}
Fix \(u\).
\begin{enumerate}[label=(\alph*),leftmargin=2em,itemsep=1pt]
\item \((\mathrm{MC}_\eta)\) for all balls in the frame \(u\), at all centres in \(U\) and all
radii, is equivalent to \(\Glin_{uu}=(2\pi/\eta)\dT_{uu}\) in \(U\).
\item Let \(\dT\) be stationary in the frame \(u\), and let \(\dT'\) be its time reverse,
\(T'_{00}=T_{00}\), \(T'_{0i}=-T_{0i}\), \(T'_{ij}=T_{ij}\). Then \(\dT'\) is conserved,
traceless and stationary, and the modular energies of all balls at rest in \(u\) are the
same for both sources. For a ball boosted with rapidity \(\beta\) in the direction \(n\)
they differ by \(2\pi\int_Bw_R\,2\sinh(2\beta)\,n^iT_{0i}\).
\item Consequently, if \(h\) satisfies \(\Glin=(2\pi/\eta)\dT\), the pair \((\dT',h)\)
satisfies \((\mathrm{MC}_\eta)\) for every ball at rest in \(u\), although
\(\Glin_{0i}=-(2\pi/\eta)T'_{0i}\).
\end{enumerate}
\end{proposition}

\begin{proof}
(a) is Theorem~\ref{thm:equivalence} with a single frame, where the leading term gives only
the \(uu\) component. (b) Stationarity and conservation give
\(\partial_iT_{0i}=-\partial_tT_{00}=0\) and \(\partial_jT_{ij}=-\partial_tT_{0i}=0\), which
are the conservation equations of \(\dT'\); the trace is unchanged. Balls at rest in \(u\)
read only \(T_{00}\). For \(u=(\cosh\beta,\sinh\beta\,n)\),
\(T_{uu}-T'_{uu}=4\cosh\beta\sinh\beta\,n^iT_{0i}\). (c) follows from (a) and
(b).
\end{proof}

A body rotating one way and the same body rotating the other way have identical ball data
in their common rest frame, and opposite frame dragging. The momentum constraint and the
spatial components of the field equation, which carry the dynamics, are the part of the
geometry that only comparison of rest frames reveals. In this precise sense the
multi-directional character of the data is where the gravitational dynamics sits.

\subsection{Several rest frames and their stability}\label{subsec:stability}

By Corollary~\ref{cor:recovery} any open set of rest frames suffices in exact arithmetic.
With noise and finite balls the question is quantitative.

\begin{proposition}[Rapidity sets the precision]\label{prop:stability}
Let the observer use its rest frame and boosts of rapidity \(\beta\) in a fixed set of
directions for which the leading-order readings determine a traceless symmetric tensor,
and let each normalized reading \(y_a=\dT_{u_au_a}(x_0)+e_a\) carry an independent error of
variance \(\sigma^2\). As \(\beta\to0\) the least-squares estimate has root-mean-square
error of order \(\sigma\) in \(T_{00}\), \(\sigma/\beta\) in \(T_{0i}\) and \(\sigma/\beta^2\)
in \(T_{ij}\).
\end{proposition}

\begin{proof}
With \(u_a=(1+\frac{\beta^2}2,\beta n_a)+O(\beta^3)\),
\(y_a=T_{00}+2\beta\,n_a^iT_{0i}+\beta^2(T_{00}+n_a^in_a^jT_{ij})+O(\beta^3)+e_a\). The columns
of the design matrix for \((T_{00},T_{0i},T_{ij})\) scale as \(1,\beta,\beta^2\), and for a
direction set whose leading blocks are nondegenerate the inverse scales as
\(1,\beta^{-1},\beta^{-2}\).
\end{proof}

Finite balls add a bias. A reading at radius \(R\) multiplies each Fourier component by
\(f(R|k_u|)\) (Proposition~\ref{prop:series}), and \(|k_u|\) grows like \(e^\beta\) for
large boosts. For the stresses the bias then grows like \(R^2e^{2\beta}|k|^2\) while the
noise falls like \(\sigma e^{-2\beta}\), and the best error obtainable with one radius scales
as \(\sqrt\sigma\,R\,|k|\). Table~\ref{tab:stability} shows both regimes for one mode in
\(d=3\), with thirteen frames (the rest frame and boosts along the six coordinate
half-axes and six face diagonals). A finer scan gives an optimal error proportional to
\(R^{1.03}\) when noise dominates and to \(\sigma^{0.48}\) between \(\sigma=10^{-4}\) and
\(10^{-2}\) at \(R=0.02L\); at very small \(\sigma\) the rest-frame bias, proportional to
\(R^2\), dominates. Since the transfer function is known exactly, readings at
several radii can remove the finite-size bias. What cannot be removed is the amplification
of noise by a small range of rapidities.

\begin{table}[ht]
\centering
\small
\begin{tabular}{@{}rcccccc@{}}
\toprule
& \multicolumn{3}{c}{\(\sigma=10^{-4}\)} & \multicolumn{3}{c}{\(\sigma=10^{-2}\)}\\
\cmidrule(lr){2-4}\cmidrule(l){5-7}
\(\beta\) & \(T_{00}\) & \(T_{0i}\) & \(T_{ij}\) & \(T_{00}\) & \(T_{0i}\) & \(T_{ij}\)\\
\midrule
0.05 & \(3.2\times10^{-5}\) & \(1.1\times10^{-3}\) & \(8.8\times10^{-2}\) & \(3.3\times10^{-3}\) & \(1.0\times10^{-1}\) & \(9.3\)\\
0.2 & \(3.1\times10^{-5}\) & \(2.7\times10^{-4}\) & \(5.8\times10^{-3}\) & \(3.1\times10^{-3}\) & \(2.7\times10^{-2}\) & \(5.8\times10^{-1}\)\\
0.7 & \(3.0\times10^{-5}\) & \(7.6\times10^{-5}\) & \(4.3\times10^{-4}\) & \(1.9\times10^{-3}\) & \(5.3\times10^{-3}\) & \(3.8\times10^{-2}\)\\
1.0 & \(5.3\times10^{-5}\) & \(9.5\times10^{-5}\) & \(3.4\times10^{-4}\) & \(1.1\times10^{-3}\) & \(3.1\times10^{-3}\) & \(1.7\times10^{-2}\)\\
2.0 & \(4.5\times10^{-4}\) & \(7.3\times10^{-4}\) & \(2.2\times10^{-3}\) & \(4.6\times10^{-4}\) & \(8.2\times10^{-4}\) & \(2.7\times10^{-3}\)\\
\bottomrule
\end{tabular}
\caption{Root-mean-square reconstruction error of a unit-norm traceless conserved source
\(\mathrm{Re}(\varepsilon_{\mu\nu}e^{ik\cdot x})\), \(k=(1,0.6,0.3,-0.2)/L\), from leading-order
readings of balls of radius \(R=0.05L\) in thirteen frames of rapidity \(\beta\), 200 noise
draws.}
\label{tab:stability}
\end{table}

\subsection{Free radiation}\label{subsec:radiation}

By Corollary~\ref{cor:weyl}, \((\mathrm{MC}_\eta)\) fixes \(h\) only up to solutions of
\(\Glin=0\). For \(d\ge3\) these include gauge modes and free gravitational waves. This is exactly the
freedom that the Einstein equation itself leaves. No ball functional registers a free wave
at linear order, in any frame and at any radius. At second order its quadratic geometry enters
$G^{(2)}[h,h]$. Moving that term to an effective source is an additional
convention, not an extra matter stress tensor. Shape-sensitive regions can also
probe tidal data: a deformed region couples to the tidal part of the curvature. In \(2+1\)
dimensions the Riemann tensor is algebraically determined by the Ricci tensor, so
\(\Glin=0\) implies a locally pure-gauge perturbation on a contractible domain.
Global identifications and boundary conditions can still distinguish metrics.

\subsection{A locally prepared source and its boundary datum}\label{subsec:prepared}

The source witness of Theorem~\ref{thm:uniqueness} can be prepared by a local operation.
For a free scalar one must also specify which part of its modular energy is being read.
We use the scalar boundary prescription of Casini, Mazzitelli and
Test\'e~\cite[Eq.~(58)]{CasiniMazzitelliTeste2015}, with the boundary-condition
derivation of Herzog and Nishioka~\cite[Secs.~1--2]{HerzogNishioka2016}.
The prescription agrees with thermal numerical calculations and R\'enyi-entropy
checks~\cite{Herzog2014,LLPS2015}. For the conformal improved stress tensor it gives
\begin{equation}\label{eq:scalaredge}
H_B=K_B+Q_B,\qquad
Q_B=2\pi\xi\int_{\partial B}\sum_a:\!\phi_a^2\!:\,dA,\qquad
\xi=\frac{d-1}{4d}.
\end{equation}
There is one real component for a real scalar and two for a complex scalar. The formula
uses vacuum-subtracted expectations. A surface expectation is defined as the limit of
smooth collar smearings of the normal-ordered field square, whose weights converge to
surface measure. No sharp boundary operator on arbitrary states is assumed.

\begin{proposition}[Local preparation supplies a nonzero spectral witness]\label{prop:prepared}
In the free massless complex scalar vacuum, let \(f\) be a real smooth compactly
supported spacetime function with nonzero one-particle class. Put
\(W_\theta=\exp[i\theta\phi_1(f)]\), \(\theta\ne0\), and
\(\varphi_f=\Delta_{\mathrm{PJ}}f\), with the sign convention fixed so that
\(\omega_{W_\theta}(\phi_1(x))=\theta\varphi_f(x)\).
Then
\begin{equation}\label{eq:preparedsource}
\dT_{\mu\nu}=\theta^2\left[
\partial_\mu\varphi_f\partial_\nu\varphi_f
-\frac12\eta_{\mu\nu}(\partial\varphi_f)^2
+\xi(\eta_{\mu\nu}\Box-\partial_\mu\partial_\nu)\varphi_f^2\right].
\end{equation}
It is smooth, conserved and traceless, has finite strictly positive total spatial energy,
and has a finite tensor-valued Radon Fourier measure with nonzero mass away from \(k=0\).
All polynomial moments of that measure are finite. Its boundary-energy variation is
\begin{equation}\label{eq:preparededge}
\delta\ev{Q_B}=2\pi\xi\theta^2\int_{\partial B}\varphi_f^2\,dA.
\end{equation}
If \(\operatorname{supp}f\) is compactly contained in \(D(B)\), this term vanishes for
that ball. It need not vanish for the same state on other balls.
\end{proposition}

The proof is in Appendix~\ref{app:prop:prepared}.


\begin{example}[A common state, source and bulk response]\label{ex:coherent}
Take the prepared state above and \(\kappa=2\pi/\eta\). A smooth linearized response
\(\Glin[h]=\kappa\dT\) exists with constraint-compatible initial data. For example,
write \(\bar h_{\mu\nu}=h_{\mu\nu}-\eta_{\mu\nu}h^\alpha{}_\alpha/2\) and solve
\(\Box\bar h_{\mu\nu}=-2\kappa\dT_{\mu\nu}\).
At a fixed time \(t_0\), choose smooth Poisson solutions
\[
\Delta a=-2\kappa\dT_{00}(t_0,\cdot),\qquad
\Delta b_j=2\kappa\dT_{0j}(t_0,\cdot).
\]
Set \(\bar h_{00}=a\), all other initial components zero,
\(\partial_t\bar h_{0\nu}=0\), and
\(\partial_t\bar h_{ij}=\partial_i b_j+\partial_j b_i-\delta_{ij}\partial_\ell b_\ell\).
The Lorenz constraint \(C_\nu=\partial^\mu\bar h_{\mu\nu}\) and its first time derivative
vanish initially. Conservation gives \(\Box C_\nu=0\), so they vanish throughout the
linear evolution and \(\Glin[h]=\kappa\dT\).
Compact smooth sources on a slice admit these Poisson solutions; in \(d=2\) logarithmic
behaviour at spatial infinity is allowed. The construction is smooth and tempered, and
is perturbatively small on each fixed compact domain as \(\theta\to0\).

For this pair Theorem~\ref{thm:equivalence} gives, for every ball,
\begin{equation}\label{eq:edgeMC}
\eta\,\delta F_B[h]
+\delta\ev{H_B}-\delta\ev{Q_B}=0.
\end{equation}
Both \(\widehat{\dT}\) and \(\widehat{\Glin[h]}=\kappa\widehat{\dT}\) meet the Radon
hypotheses. Hence one radius, all centres and an open frame family identify the
canonical intrinsic response from these edge-subtracted readings.
The matter state, its physical source and the response law coexist here; the Einstein
coupling and the constraint-compatible gravitational data were supplied, not generated
by the Weyl preparation.
\end{example}

The subtraction has operational content. It requires access to calibrated collar
readings of \(\sum_a:\!\phi_a^2\!:\), in addition to the modular energy or the bulk stress
readings. Their limiting expectations are finite in this coherent family, but no finite
sample, resolution or preparation cost for the collar limit has been derived.
The total-energy formula counts the prepared radiation, not the switching,
phase-reference or reset costs of the source apparatus.
If the bulk stress charge is read directly, the subtraction is unnecessary; calling that
charge the raw physical \(H_B\) would still be incorrect.

\begin{corollary}[Raw scalar energy cannot obey the bulk law on all small balls]\label{cor:edgeobstruction}
At every centre \(x_0\) with \(\varphi_f(x_0)\ne0\), and in any one fixed rest frame,
no smooth metric perturbation satisfies
\(\eta\,\delta F_B[h]+\delta\ev{H_B}=0\) for all sufficiently small radii.
Thus the obstruction is visible at each point of a nonempty open set for every
nonzero coherent preparation.
\end{corollary}

\begin{proof}
At a point with \(\varphi_f(x_0)\ne0\),
\[
\delta\ev{Q_B}
=2\pi\xi\theta^2\Omega_{d-1}R^{d-1}\varphi_f(x_0)^2+O(R^{d+1}).
\]
The bulk charge and the canonical intrinsic geometric variation are both
\(O(R^{d+1})\) for smooth source and metric fields. They cannot cancel the nonzero
\(R^{d-1}\) term at that centre, even without varying the rest frame.
\end{proof}

\begin{remark}[The same curvature coupling fixes the surface subtraction]\label{rem:waldedge}
The boundary coefficient has a gravitational counterpart, not merely a readout
calibration. For the nonminimally coupled scalar Lagrangian
\[
\mathcal L=\left(\frac{1}{16\pi G}-\frac{\xi}{2}\Phi^2\right)\mathcal R
-\frac12\sum_a(\nabla\phi_a)^2,\qquad \Phi^2=\sum_a\phi_a^2,
\]
the binormal curvature derivative gives the Wald surface density
\cite[Eqs.~(53), (93)]{IyerWald1994}
\[
S_{\mathrm W}[\Sigma]=\int_\Sigma
\left(\frac{1}{4G}-2\pi\xi\Phi^2\right)dA.
\]
The nonminimal scalar contribution and its relation to the entropy contact term
are established results; see, for example, Donnelly and Wall~\cite[Eq.~(5),
Sec.~II.A]{DonnellyWall2012}. Here the same coefficient is used in the ball-family
consistency law.
Indeed, for the curvature coefficient \(c\),
\(\partial\mathcal L/\partial R_{abcd}
=c(g^{ac}g^{bd}-g^{ad}g^{bc})/2\);
contraction with a binormal of norm \(-2\) and multiplication by \(-2\pi\)
give \(4\pi c\). In the coherent witness the vacuum-subtracted quadratic profile is
\(\delta\ev{\Phi^2}=\theta^2\varphi_f^2\). At the retained order, with
\(h=O(\theta^2)\), its surface-measure correction is \(O(\theta^4)\).
For \(\eta=1/(4G)\), the Wald-matched ball quantity, with the same supplied
volume term as \(F_B\), is consequently
\[
\mathcal B_B:=\delta S_{\mathrm W}[\partial B]
-\eta\frac{d-1}{R}\delta V(B)
=\eta\,\delta F_B-\delta\ev{Q_B}.
\]
Thus \(\mathcal B_B+\delta\ev{H_B}
=\eta\,\delta F_B+\delta\ev{K_B}\): the bulk consistency law is exactly restored.
This matches a local surface density and fixes its matter coefficient from the same
curvature coupling. It is a coefficient identity for the ball consistency law, not
a thermodynamic first law for arbitrary diamonds.\footnote{An arbitrary ball edge
is not a Killing horizon; MC, a generalized volume and a nonlinear second law do
not follow from this density match.}
The pure-area/raw-energy obstruction above is unchanged.
\end{remark}

\section{Which algebra the observer has}\label{sec:selection}

Consciousness~I and Foundation~I select an observer's algebra by noise and by a readable
seed, and the selected algebra is in general smaller than the full local
algebra~\cite{PaperI,FoundationI}. HAFF-D conjectures that gravitational phenomena are tied
to the selection of observable algebras~\cite{HAFFD}. In the ball family the question is
precise: does the geometry fixed by \((\mathrm{MC}_\eta)\) depend on which algebra is
selected?

For a bounded collection of scalar preparations and tested stress or collar
smearings, choose a containing diamond \(D_0\). The \(U(1)\) Haar expectation \(E\)
on this diamond preserves the vacuum and fixes the invariant readings.
The semigroup \(\mathcal E_t=e^{-\gamma t}\mathrm{id}+(1-e^{-\gamma t})E\) is
vacuum-stationary. Foundation~I's stationary-source and stable-seed
theorems~\cite{FoundationI} place the modular hull of any such invariant seed inside
the proper fixed algebra. The charged coherent preparation above requires a separate
phase-referenced source. A neutral quadratic seed supplies another preparation,
entirely inside the same fixed algebra.

\begin{proposition}[A neutral preparation inside a stable hull]\label{prop:neutral}
In the free complex scalar theory of Proposition~\ref{prop:prepared}, let
\(X_a=\phi_a(f)\), \(V_f=\omega(X_a^2)>0\), and
\[
Q_f=\tfrac12(X_1^2+X_2^2),\qquad U_\theta=e^{i\theta Q_f}.
\]
The strongly commuting fields \(X_1,X_2\) define a self-adjoint \(Q_f\) by their
joint spectral calculus. Both \(Q_f\) and \(U_\theta\) are \(U(1)\)-invariant and
local to any diamond containing \(\operatorname{supp}f\).
Include \(\arctan Q_f\), together with the desired bounded stress and collar
readings, in the stable seed above. Its hull contains \(U_\theta\) for every
\(\theta\), while remaining inside the proper fixed algebra.

Put \(v_f=\varphi_f\), and define
\(H_f(x)=\tfrac12\omega(\{X_a,\phi_a(x)\})\) for either fixed component
\(a=1,2\), with no sum over \(a\); internal symmetry makes the two functions equal.
Let \(t[v]\) denote the quadratic tensor in brackets in~\eqref{eq:preparedsource},
and polarize it as \(t[v,z]=(t[v+z]-t[v]-t[z])/2\). The exact source is
\begin{equation}\label{eq:neutralsource}
\dT_{\mu\nu}=4\theta\,t_{\mu\nu}[v_f,H_f]
                 +2\theta^2 V_f\,t_{\mu\nu}[v_f].
\end{equation}
For every \(\theta\ne0\), it is smooth, conserved and traceless, has compact support
on each spatial slice and has a finite Radon Fourier measure with all polynomial
moments finite and nonzero weight away from zero. Its total energy is
\[
\int\dT_{00}\,d^dx
=\theta^2 V_f\int\big[(\partial_t v_f)^2+|\nabla v_f|^2\big]d^dx>0.
\]
Its boundary datum is
\(\delta\ev{\Phi^2}=4\theta v_fH_f+2\theta^2V_fv_f^2\).
\end{proposition}

The proof is in Appendix~\ref{app:prop:neutral}.


The constraint construction of Example~\ref{ex:coherent} gives
\(\Glin[h]=(2\pi/\eta)\dT\), so this preparation identifies the canonical response.
Its inputs are permission to apply the neutral quadratic pulse and the complete
centre and frame family. Containing the stress tensor need not supply \(Q_f\).
Nor does neutrality by itself guarantee a nonzero tangent: the local conformal
Wick square \(:\Phi^2:(g)\) can have zero first-order vacuum stress response.
The spatially smeared square used here is a different operator, and its positive
finite-amplitude energy is computed explicitly. The noise fixes the invariant
readings, although the prepared state need not be stationary.

There is a more general finite-energy route that uses a bounded observable already
in the selected algebra. It needs no particular quadratic seed. Local source
regularity remains a separate condition.

\begin{proposition}[A bounded accessible preparation with positive energy]
\label{prop:bounded-source}
Let \(H\ge0\) be the translation generator, with
\(\ker H=\mathbb C\Omega\). Let \(\mathcal N\) be a selected von Neumann algebra
for which \(\Omega\) is separating. Suppose it contains a bounded non-scalar
self-adjoint \(C\), together with \(\alpha_t(C)=e^{itH}Ce^{-itH}\) for
\(0<t<\tau\), where \(\tau>0\). For a sufficiently narrow real nonnegative
\(f\in C_c^\infty\) in that interval, \(\int f=1\), put
\[
 B=\int f(t)\alpha_t(C)dt,\qquad U_\theta=e^{i\theta B}.
\]
Then \(B\) is bounded, self-adjoint and non-scalar, \(U_\theta\in\mathcal N\),
and its vacuum excitation has finite energy for every \(\theta\). For every
\(\theta\) for which \(U_\theta\) is non-scalar, in particular
\(0<|\theta|\|B\|<\pi\),
\[
 0<E_\theta:=\langle U_\theta\Omega,H U_\theta\Omega\rangle
 \le|\theta|\|[H,B]\|.
\]
With \(K=[H,B]\), \(\psi=(C-\langle\Omega,C\Omega\rangle)\Omega\) and the
spectral measure \(\mu_\psi\) of \(H\), the quadratic coefficient is
\begin{align}
 E_2&=\langle B\Omega,HB\Omega\rangle
 =\tfrac12\langle\Omega,[B,[H,B]]\Omega\rangle\nonumber\\
 &=\int_{(0,\infty)}\lambda|\widehat f(\lambda)|^2\,d\mu_\psi(\lambda)>0,
 \qquad E_2\le\|B\|\|K\|, \label{eq:bounded-energy}
\end{align}
where \(\widehat f(\lambda)=\int f(t)e^{it\lambda}dt\). An explicit remainder is
\[
 |E_\theta-\theta^2E_2|
 \le2|\theta|^3\|B\|^2\|K\|
       +\tfrac12\theta^4\|B\|^3\|K\|.
\]
If the state's vacuum-subtracted energy density is integrable on a time slice
and its integral is \(E_\theta\), this gives a nonzero integrated source.
Implementation permission and its cost are additional inputs.
\end{proposition}
The proof is in Appendix~\ref{app:prop:bounded-source}.


The smoothed operator has bounded iterated commutators
\(\operatorname{ad}_H^m(B)=i^m\int f^{(m)}(t)\alpha_t(C)dt\).
Repeated use of the product rule in the exponential series therefore gives
\(U_\theta\Omega\in\bigcap_m\operatorname{Dom}H^m\); all energy moments are
finite. If \(\mu_\psi\) has no atoms, narrowness is unnecessary: the nonzero
entire function \(\widehat f\) has only isolated zeros, which cannot carry
that spectral measure. With atoms, a broad smearing can cancel an excitation
exactly, so this refinement requires the stated spectral condition.

\begin{corollary}[A common pulse inside a shared hull]\label{cor:shared-pulse}
Let several selected algebras \(\mathcal N_j\) act in one vacuum representation
with one positive generator \(H\), \(\ker H=\mathbb C\Omega\), and \(\Omega\)
separating for an ambient algebra containing them. Suppose each is invariant
under \(\operatorname{Ad}e^{itH}\) for \(t\ge0\). If their intersection
\(\mathcal N_\cap\) is non-scalar, it contains one common bounded pulse
\(U_\theta\) with finite strictly positive \(H\)-energy.
\end{corollary}
\begin{proof}
The intersection is a von Neumann algebra, inherits separation and forward
translation invariance, and contains a bounded non-scalar self-adjoint \(C\).
Apply Proposition~\ref{prop:bounded-source} in \(\mathcal N_\cap\).
\end{proof}

\begin{proposition}[Finite algebras cannot carry this vacuum permission]
\label{prop:finite-vacuum-permission}
Let $\mathcal N$ be finite-dimensional, let $H\ge0$ have
$\ker H=\mathbb C\Omega$, and suppose $\Omega$ is separating for
$\mathcal N$. If $\operatorname{Ad}e^{itH}(\mathcal N)\subseteq
\mathcal N$ for every $t\ge0$, then $\mathcal N=\mathbb C1$.
The ambient Hilbert space need not be finite-dimensional.
\end{proposition}
\begin{proof}
Injectivity and finite dimension turn each inclusion into equality, so
the automorphisms extend to negative times. Continuity fixes every
minimal central projection. On each matrix block choose a self-adjoint
generator $K$ and matrix units $E_{ab}$ in its eigenbasis:
$\alpha_t(E_{ab})=e^{it(k_a-k_b)}E_{ab}$.
Since $H\Omega=0$ and separation makes every nonzero $E_{ab}\Omega$
nonzero, this vector is an $H$-eigenvector with eigenvalue $k_a-k_b$.
Positivity applied to $E_{ab}$ and $E_{ba}$ makes all differences zero.
Thus the automorphisms fix $\mathcal N$, and $C\Omega$ lies in the
one-dimensional kernel of $H$ for every $C\in\mathcal N$. Separation
then gives $C=c1$.
\end{proof}

The common pulse of Corollary~\ref{cor:shared-pulse} consequently uses
an infinite field algebra when its full forward permission is exact.
A finite observer can implement a finite-duration pulse or an explicitly
controlled approximation; its controller Hamiltonian cannot simply be
identified with the vacuum translation generator. The bounded-pulse
construction requires only its stated finite interval and is unaffected.

Foundation~I's \(\Gamma_3\)-hulls have precisely this forward translation
closure~\cite[Theorem on Cut inheritance]{FoundationI}. Thus a non-scalar
shared intersection of those hulls, with the common representation hypotheses
above, supplies a common energy-orientation test without adding \(Q_f\).
The positive generator must be the same for the readers: in a null Cut model it
is null momentum, not automatically the time-translation Hamiltonian of this
higher-dimensional theory. A merely modular-stable \(\Gamma_2\)-hull does
not inherit translation permission. Realizing the pulse and reading its energy
still require controls and an energy ledger.

There is a common physical-$H$ instance in $3+1$ dimensions.
For two independent free conformal complex scalar nets, take the
future-cone algebra $\mathcal A(V_+)$ and the two selected nets
$\mathcal A_1^{U(1)}\bar\otimes\mathcal A_2$ and
$\mathcal A_1\bar\otimes\mathcal A_2^{U(1)}$.
Time translations move $V_+$ into itself, the vacuum is separating
there, and the physical $H_1+H_2$ has unique zero-energy vacuum
\cite{HislopLongo1982,BGL1993}. Each fixed algebra is taken on its own
vacuum-cyclic subspace. The same physical Hamiltonian restricts to it
because the Haar projection commutes with translations.
Their intersection contains the neutral preparation of
Proposition~\ref{prop:neutral}, the total stress and the invariant
collar readings at each tested diamond. Thus the reference cone
algebra and the selected ball family coexist in one matter model;
they are distinct region algebras, not all operators inside one cone.
The generic bounded pulse keeps its separate source-regularity
condition. Shared implementation, calibrated collar limits and
$\mathrm{MC}_\eta$ remain supplied permissions and consistency data.

\begin{corollary}[A positive pulse supplies the spectral witness]
\label{cor:pulse-witness}
Suppose a pulse from Proposition~\ref{prop:bounded-source} has a smooth source
\(\dT\) in the tempered Radon Fourier class, and its vacuum-subtracted energy
density is integrable on a time slice with integral \(E_\theta>0\).
Then its source Fourier measure has nonzero weight away from \(k=0\).
If the response also belongs to that class and the full centre and open-frame
consistency hypotheses of Theorem~\ref{thm:uniqueness} hold, this source identifies
the canonical response up to normalization.
\end{corollary}
\begin{proof}
A Radon measure supported only at zero is a multiple of \(\delta_0\);
the source would be constant. Its energy density is then constant on the
infinite spatial slice. Integrability forces that component to vanish,
contradicting \(E_\theta>0\). The identification follows from
Theorem~\ref{thm:uniqueness}.
\end{proof}

This witness uses finite amplitude, not necessarily a nonzero first-order
stress tangent. In fact
\(\partial_\theta E_\theta|_0=i\langle\Omega,[H,B]\Omega\rangle=0\)
for every such vacuum pulse. A local tangent depends separately on
\(-2\operatorname{Im}\langle\Omega,T(x)B\Omega\rangle\).
Proposition~\ref{prop:neutral} computes both the regularity and a source
explicitly; the general bounded construction does not supply Radon regularity
from finite energy moments alone. Neither route establishes finite-resource
identification of the entire geometric response.

\begin{proposition}[Selection is invisible to bulk modular energy]\label{prop:selection}
For each ball let \(\mathcal A_{\mathrm{sel}}(B)\subset\mathcal A(D(B))\) be a von Neumann
subalgebra with a normal conditional expectation \(E_B\) that preserves the vacuum state,
and suppose \(\mathcal A_{\mathrm{sel}}(B)\) contains the bounded functions of stress-tensor
smearings in \(D(B)\); fixed-point algebras of a compact internal symmetry are examples.
Then the modular group of \((\mathcal A_{\mathrm{sel}}(B),\omega)\) is the restriction of that
of \(\mathcal A(D(B))\), the quantity \(\delta\ev{K_B}=2\pi\int_Bw_R\,\psi(T_{uu})\) is
determined by the restriction of the state \(\psi\) to \(\mathcal A_{\mathrm{sel}}(B)\), and
\((\mathrm{MC}_\eta)\) is the same condition for every such choice. Every such selected
algebra therefore fixes the same matter-sourced Einstein tensor. For \(d\ge3\), the
metric additionally requires common initial or boundary data and is defined modulo
gauge; for \(d=2\), it is locally fixed up to gauge on contractible domains.
If the physical energy has an extra boundary term, the bulk datum must be read directly
or that term must also be accessible and subtracted.
\end{proposition}

\begin{proof}
The first statement is Takesaki's theorem on conditional
expectations~\cite{Takesaki1972}. The second holds because the stress-tensor smearings are
affiliated with \(\mathcal A_{\mathrm{sel}}(B)\). The last follows from
Theorem~\ref{thm:equivalence}.
\end{proof}

Relative entropy behaves differently, as Foundation~I found for current excitations of the
SU(2) level-one net~\cite{FoundationI}. The same argument works in any dimension.


\begin{proposition}[Selection changes entropy]\label{prop:entropy}
Let \(\mathcal A\) be the net of a free massless conformally coupled complex scalar
\(\phi=(\phi_1+i\phi_2)/\sqrt2\) with its internal \(U(1)\), let
\(\mathcal A_{\mathrm{sel}}(B)\) be the \(U(1)\)-fixed points of \(\mathcal A(D(B))\) and \(E_B\)
the average over the group, and let \(W=e^{i\theta\phi_1(f)}\), \(\theta\ne0\), where \(f\) is real, supported
in \(D(B)\), and has nonzero one-particle class. A nonzero test function alone
would not suffice, since \(f=\Box g\) can represent the zero field.
\begin{enumerate}[label=(\alph*),leftmargin=2em,itemsep=1pt]
\item \(S_{\mathcal A_{\mathrm{sel}}(B)}(\omega_W\Vert\omega)<S_{\mathcal A(D(B))}(\omega_W\Vert\omega)
=\delta\ev{K_B}_{\omega_W}\).
\item As \(\theta\to0\), the modular energy is quadratic in \(\theta\), and the
restricted state has vanishing first variation. At a faithful finite-dimensional regulator,
write \(\rho^{\mathrm{inv}}_\theta=\rho_0+\theta^2\chi+O(\theta^4)\). Then
\[
S(\rho^{\mathrm{inv}}_\theta\Vert\rho_0)
=\frac{\theta^4}{2}g_{\mathrm{BKM},\rho_0}(\chi,\chi)+O(\theta^6).
\]
The continuum quartic law additionally requires a finite BKM tangent norm and
uniform control of the entropy remainder; smooth field smearing and faithfulness
alone do not establish these conditions.
\end{enumerate}
\end{proposition}

The proof is in Appendix~\ref{app:prop:entropy}.


A local coherent displacement preserves the regulated entropy of the full algebra.
Its modular energy equals its relative entropy. The invariant algebra strictly loses
relative entropy while retaining the stress charge; when the quartic expansion is
controlled, its entropy increase exhausts that charge to order \(\theta^2\).

Appendix~\ref{app:bkm} computes the invariant mode coefficient and separates the
exponential same-charge contribution from neutral-pair terms. A bounded BKM tangent
norm is only one requirement: a continuum quartic entropy law still needs uniform
control of the remainder as the regulators are removed. The strict entropy inequality
above is independent of that expansion.

\begin{conjecture}[Selection enters through second-order quantities or through the source]\label{conj:selection}
In the conformal ball setting, a change of the selected algebra at a fixed state changes the
geometry reconstructed from ball modular data only through quantities of second order in
the state perturbation, such as relative entropy, Fisher information and canonical energy,
or through physical operations that change the stress tensor.
\end{conjecture}

\emph{Mechanism.} The first-order datum is a function of the restriction of the state to
stress-tensor observables, which every relevant algebra contains
(Proposition~\ref{prop:selection}). What sees more of the state is relative entropy, which
depends on the algebra by monotonicity, and in holographic theories its second-order part
for balls is the canonical energy that carries the second-order gravitational
equations~\cite{LashkariVanRaamsdonk2016,FHHPRV2017}.

\emph{Support.} Propositions~\ref{prop:selection} and~\ref{prop:entropy}; Foundation~I's
strict inequality in \(1+1\) dimensions~\cite{FoundationI}; the result of HAFF-D that a
record alone does not move a gravitational source, while a preparation
does~\cite{HAFFD}.

\emph{Strongest objection.} At second order the world still has one metric. The relative
entropy of a smaller algebra measures what its observer cannot distinguish, so a
second-order condition built from it may describe a coarser reconstruction rather than a
different geometry. If so, the conjecture is a statement about the reach of an observer's
reconstruction, which an observer-relative account of geometry has to quantify anyway,
and not about a second spacetime.

\emph{Test.} Compute \(S_{\mathcal A(D(B))}-S_{\mathcal A_{\mathrm{sel}}(B)}\) for the charged
coherent excitation of Proposition~\ref{prop:entropy} through order \(\theta^4\), and compare
the second-order consistency conditions built from the two relative entropies.

\section{Second-order geometry and a source that follows it}
\label{sec:secondorder}

The linear identity holds at every radius. At second order two of its
inputs change: the intrinsic curvature need not equal twice the normal
Einstein tensor, and the flat CHM charge need not be the modular Hamiltonian
of the curved theory. We can resolve the first issue and identify the source
terms needed for the second. This does not supply a nonlinear modular law.

\subsection{A geometric two-jet consistency law}

Put $c_0=V_d/(d+2)$. Let $g_\epsilon$ be a smooth family of Lorentzian
metrics with $g_0=\eta$. At $p$, choose smooth spacelike hypersurfaces
$\Sigma_{\epsilon,u}$ whose flat member is the hyperplane normal to $u$.
Let $n_\epsilon$ be the future unit normal, $q_\epsilon$ its induced
metric, and $K_\epsilon$ its second fundamental form. The family, including
its centre and normal prescription, is part of the experiment.
For the intrinsic geodesic ball of proper radius $R$, define
\begin{align}
 \mathfrak F_{R,u}(g_\epsilon)
 &=A_\epsilon-\frac{d-1}{R}V_\epsilon
       -\frac{\Omega_{d-1}}d R^{d-1}
       -\frac{c_0}2R^{d+1}\mathcal Q_\epsilon(p),
       \label{eq:secondF}\\
 \mathcal Q_\epsilon&=(\operatorname{tr}_{q_\epsilon}K_\epsilon)^2
                 -K_{\epsilon,ij}K_\epsilon^{ij}.\nonumber
\end{align}
The last term is quadratic about a flat slice. It does not change any of
the linear results, but it changes the second-order claim. This is a
geometric correction dictated by the Gauss constraint, not an entropy
or modular correction derived from a quantum state.
The proper radius $R$ and its multiplier $(d-1)/R$ are held fixed as
$\epsilon$ varies. A coordinate sphere of unchanged coordinate radius
is a different experiment: its geodesic radius and boundary measure
change. The radial check below explicitly inverts that distance change.

For a smooth symmetric source $T_\epsilon$, define a \emph{stress reading}
\begin{equation}
 \mathfrak K_{R,u}(T_\epsilon,g_\epsilon)
 =2\pi\int_{B_\epsilon}\frac{R^2-r^2}{2R}
                T_{\epsilon,\mu\nu}n_\epsilon^\mu n_\epsilon^\nu\,dV_\epsilon.
 \label{eq:secondK}
\end{equation}
Here $r$ is intrinsic geodesic distance, and the normal field on the
slice is used inside the integral. In flat spacetime this is the bulk
CHM charge. On a general curved metric it is only the specified local
stress reading. No assertion of geometric modular flow on the curved
diamond is being made.

\begin{proposition}[Leading nonlinear ball identity]
\label{prop:nonlinear-leading}
On a compact family with bounded smooth geometric and source derivatives,
spacelike slices and radii below their normal injectivity radius,
\begin{align}
 \mathfrak F_{R,u}&=-c_0R^{d+1}G_{\mu\nu}(g_\epsilon)(p)
                         n_\epsilon^\mu n_\epsilon^\nu+O(R^{d+3}),\\
 \mathfrak K_{R,u}&=2\pi c_0R^{d+1}T_{\epsilon,\mu\nu}(p)
                         n_\epsilon^\mu n_\epsilon^\nu+O(R^{d+3}).
 \label{eq:nonlinear-leading}
\end{align}
The remainders and their first two $\epsilon$ derivatives are uniform
when the corresponding mixed derivatives are bounded.
\end{proposition}
\begin{proof}
In intrinsic normal coordinates the volume density is
$1-\operatorname{Ric}_{ij}x^ix^j/6+O(|x|^3)$. Angular averaging removes
the odd terms and gives
\[
 V=V_dR^d\left(1-\frac{{}^{(d)}\!\mathcal R}{6(d+2)}R^2\right)
       +O(R^{d+4}),\qquad
 A=\Omega_{d-1}R^{d-1}\left(1-\frac{{}^{(d)}\!\mathcal R}{6d}R^2\right)
       +O(R^{d+3}).
\]
Thus the area--volume difference minus its flat value is
$-c_0R^{d+1}{}^{(d)}\!\mathcal R/2+O(R^{d+3})$.
The exact Gauss equation
$ {}^{(d)}\!\mathcal R+\mathcal Q=2G_{nn}$ proves the first statement.
Taylor expansion of the smooth scalar $T_{nn}$ and the volume density
proves the second; its constant integral is $c_0R^{d+1}$ and its odd
terms average to zero. Differentiating these local expansions establishes
the stated uniformity. Sufficient bounded mixed derivatives are assumed,
rather than inferred from finite energy alone.
\end{proof}

Write, with no factorial in the second coefficient,
\begin{align}
 g_\epsilon&=\eta+\epsilon h+\epsilon^2j+O(\epsilon^3),
 & T_\epsilon&=\epsilon t_1+\epsilon^2t_2+O(\epsilon^3),\\
 G(g_\epsilon)&=\epsilon\Glin[h]
       +\epsilon^2\bigl(\Glin[j]+G^{(2)}[h,h]\bigr)+O(\epsilon^3).
 \label{eq:twojets}
\end{align}
$G^{(2)}$ is the quadratic coefficient of the covariant Einstein tensor
in this fixed smooth identification of the metrics. Coefficients of the
source use that same identification and the same renormalization convention.
Appendix~\ref{app:twojet} gives its explicit coordinate formula and the
source variation used below.

\begin{theorem}[Second-order small-ball equivalence]
\label{thm:secondorder}
Suppose $\Glin[h]=(2\pi/\eta)t_1$ on an open domain and the two-jet
families above have the regularity of Proposition~\ref{prop:nonlinear-leading}.
Then the following are equivalent:
\begin{enumerate}[label=(\alph*),leftmargin=2em,itemsep=1pt]
\item At every $p$ and for $u$ in a nonempty open frame set,
\[
 \lim_{R\to0}R^{-d-1}
 [\epsilon^2]\bigl(\eta\mathfrak F_{R,u}
                         +\mathfrak K_{R,u}\bigr)=0.
\]
\item The second-order tensor equation is
\[
 \Glin[j]+G^{(2)}[h,h]=\frac{2\pi}{\eta}t_2.
\]
\end{enumerate}
Here $[\epsilon^2]$ denotes the quadratic Taylor coefficient. For one
frame only the corresponding normal constraint follows.
\end{theorem}
\begin{proof}
Set
\[
 E(g_\epsilon)=G(g_\epsilon)-(2\pi/\eta)T_\epsilon.
\]
Proposition~\ref{prop:nonlinear-leading} gives the leading residual
$-\eta c_0R^{d+1}E_{n_\epsilon n_\epsilon}$.
Its zeroth and first coefficients vanish. Consequently changes of the
normal at order $\epsilon$ multiply the vanishing first residual and do
not alter the quadratic coefficient. That coefficient is
$-\eta c_0(\Glin[j]+G^{(2)}[h,h]-(2\pi/\eta)t_2)_{uu}$.
Tensor recovery from an open frame set proves both directions.
\end{proof}

This is a second-order conditional law in the same spacetime, not a
holographic result. Unlike Theorem~\ref{thm:equivalence}, it gives no exact
finite-radius identity. Terms quadratic in curvature, its derivatives,
the ball prescription and curved modular flow must be accounted for there.
Nor does it classify all quadratic response functionals: adding terms
whose small-ball leading coefficient vanishes leaves this test unchanged.

The extrinsic correction can instead be avoided by a particular physical
region prescription. The local exponential image at $p$ of
$n_\epsilon^\perp$ is a spacelike geodesic hypersurface with
$K_\epsilon(p)=0$. On that family the uncorrected area--volume difference
has the same leading law. It is not the same region as a ball on an
arbitrary fixed laboratory slice. Regions, normals and preparations cannot
be changed while interpreting their readings as the same experiment.

\begin{example}[A bent slice in flat spacetime]
\label{ex:bent-slice}
Keep the ambient metric exactly Minkowski and the source zero. Near the
origin choose $t=\epsilon\alpha |x|^2/2$. Its normal there is $\partial_t$,
and $K_{ij}=\pm\epsilon\alpha\delta_{ij}+O(\epsilon^2)$, so
$\mathcal Q=\epsilon^2d(d-1)\alpha^2+O(\epsilon^3)$ and
$ {}^{(d)}\!\mathcal R=-\mathcal Q$. The raw ball difference is
\[
 [\epsilon^2]\left(F_R-\frac{\Omega_{d-1}}dR^{d-1}\right)
 =\frac{\Omega_{d-1}(d-1)}{2(d+2)}\alpha^2 R^{d+1},
\]
although $G=0$. To check the coefficient directly, its radial metric is
$(1-\epsilon^2\alpha^2s^2)ds^2+s^2d\Omega^2$; invert its proper-radius
integral and expand its area and volume. The correction in
Eq.~\eqref{eq:secondF} cancels this term exactly at quadratic order.
Without a slice prescription or this correction, a nonlinear extension
would assign a source to flat spacetime merely by bending a slice.
\end{example}

\subsection{Why the flat prepared source cannot simply be reused}

Even if a quantum stress expectation is smooth, a second-order response
must use the total source in its perturbed metric. With a common covariant
renormalization prescription and no gravitational anomaly, its Ward law
is $\nabla^\mu T_{\mu\nu}=0$. Expanding that identity gives
\begin{equation}
 \partial^\mu t_{2,\mu\nu}=-\mathcal D_h[t_1]_\nu,
 \label{eq:secondward}
\end{equation}
where
\begin{align*}
 \mathcal D_h[t]_\nu={}&-h^{\mu\alpha}\partial_\alpha t_{\mu\nu}
 -\eta^{\mu\alpha}\Gamma^{(1)\lambda}_{\alpha\mu}t_{\lambda\nu}
 -\eta^{\mu\alpha}\Gamma^{(1)\lambda}_{\alpha\nu}t_{\mu\lambda},\\
 \Gamma^{(1)\lambda}_{\alpha\beta}={}&\tfrac12\eta^{\lambda\rho}
 (\partial_\alpha h_{\beta\rho}+\partial_\beta h_{\alpha\rho}
                                  -\partial_\rho h_{\alpha\beta}).
\end{align*}
Thus keeping only a conserved flat $t_1$ is generally not a second-order
source: its curved divergence is $\epsilon^2\mathcal D_h[t_1]$.
The expanded Bianchi identity demands precisely this correction.

The coherent preparation permits an explicit classical calculation of
the missing terms. Take the same nonminimal scalar coupling as in
Remark~\ref{rem:waldedge}, now varied in the curved metric. For one real
mean field $\Phi$, its stress is
\begin{equation}
 T_{\mu\nu}[g,\Phi]=\partial_\mu\Phi\partial_\nu\Phi
 -\tfrac12g_{\mu\nu}(\nabla\Phi)^2
 +\xi(g_{\mu\nu}\Box_g-\nabla_\mu\nabla_\nu)\Phi^2
 +\xi G_{\mu\nu}(g)\Phi^2.
 \label{eq:curvedscalar}
\end{equation}
Its divergence is $(\Box_g\Phi-\xi\mathcal R_g\Phi)\partial_\nu\Phi$;
the derivative of the last term cancels the curvature derivative from
the improvement. Its trace vanishes on shell for the conformal value
$\xi=(d-1)/(4d)$ at the classical level.

Set $\epsilon=\theta^2$ for the Weyl witness, and take
$g=\eta+\theta^2h_2+\theta^4h_4+\cdots$ and
$\Phi=\theta v+\theta^3v_3+\cdots$. Then
\begin{align}
 \Box v_3&=-\delta\Box_{h_2}v+\xi\mathcal R^{\mathrm{lin}}[h_2]v,
 \label{eq:scalar-return}\\
 t_2&=t[v],\qquad
 t_4=2t[v,v_3]+\delta_g t_{h_2}[v]
                       +\xi\Glin[h_2]v^2.\label{eq:fourthsource}
\end{align}
Here $t_g[v]$ comprises the first three terms of
Eq.~\eqref{eq:curvedscalar}; $\delta_g t_{h_2}[v]$ varies them with $v$
fixed, and $t[v,z]$ is their flat polarization used in
Proposition~\ref{prop:neutral}. In coordinates,
\[
 \delta\Box_h v=-h^{\mu\nu}\partial_\mu\partial_\nu v
 -(\partial_\mu h^{\mu\nu}-\tfrac12\partial^\nu h^\alpha{}_{\alpha})
                                                  \partial_\nu v.
\]
On any fixed finite causal domain with smooth $h_2$, the driven flat
wave equation for $v_3$ has smooth solutions after its Cauchy data are
specified. Taking zero correction data is one preparation prescription,
not the same as keeping the field unchanged through the perturbed evolution.
Equations~\eqref{eq:scalar-return}--\eqref{eq:fourthsource} obey
Eq.~\eqref{eq:secondward} at order $\theta^4$. The curvature response,
changed propagation and metric variation are calculable requirements,
not optional control overhead. This is an actual two-jet source
calculation within the supplied classical scalar action; it neither
derives that action nor proves a quantum source apparatus realizes it.

This calculation has nonempty second-order content. Choose smooth compact
Cauchy data with $v(p)=B\ne0$, $\partial_t v(p)=A\ne0$ and spatially
constant data near $p$. They can be represented by a local Weyl preparation
using a smooth time-slice cutoff. Then $t_{2,00}(p)=A^2/2>0$.
Choose a constraint-compatible $h_2$ as in Example~\ref{ex:coherent}
and use a linearized normal-coordinate gauge at $p$, so $h_2(p)=0$
and $\partial h_2(p)=0$. Prescribe zero $v_3$ and first time derivative
on that Cauchy slice. Tracefreeness gives
$\mathcal R^{\mathrm{lin}}[h_2](p)=0$; Eq.~\eqref{eq:scalar-return}
then gives zero second time derivative of $v_3$ at $p$ as well.
All second derivatives entering the polarized stress consequently vanish
there, and $\delta_gt_{h_2}[v](p)=0$. Therefore
\[
 t_{4,\mu\nu}(p)=\xi\kappa B^2t_{2,\mu\nu}(p),\qquad
 t_{4,00}(p)=\xi\kappa B^2A^2/2>0,
 \quad \kappa=2\pi/\eta.
\]
This is one supplied local preparation and coordinate prescription,
not a gauge-invariant attribution to one separate term on every slice.
The sum in Eq.~\eqref{eq:fourthsource} is the covariant source.

For a coherent displacement of a Hadamard reference on the \emph{same}
metric, with unchanged connected two-point function, point splitting
gives exactly this state-dependent stress difference: the added
two-point function is $\Phi(x)\Phi(y)$ and the state-independent
counterterms cancel in the difference. The total gravitational source additionally contains
the reference's renormalized stress, metric-dependent vacuum response
and the preparation/control sources. Those terms cannot be declared
zero by changing the reference independently on every ball. The scalar
surface expectation and its measure also acquire order-$\theta^4$
corrections. A Wald coefficient match remains useful, but does not turn
a generic curved diamond into a Killing horizon or supply its modular
Hamiltonian.

\subsection{A strong common-geometry extension}

\begin{conjecture}[A causal source and a common second-order geometry]
\label{conj:secondorder-common}
There is a physically preparable family in a common positive-energy
covariant matter theory, with a fixed source reference and supplied
gravitational data, for which the corrected small-ball stress consistency
of Theorem~\ref{thm:secondorder} is the local limit of the actual
modular/Wald response of its selected region family. It has the following
additional properties:
\begin{enumerate}[leftmargin=2em,itemsep=1pt]
\item The renormalized total source and its two-jet Ward identity include
the controller, old reference, support and changed propagation, rather
than only the outgoing mean pulse.
\item The geometric slice prescription and its quadratic extrinsic term
are common to the observers. Their source readings and boundary
subtractions identify one $G^{\mathrm{lin}}[j]+G^{(2)}[h,h]$ with one
normalization; changing a stress-containing algebra alone does not create
a second sourced spacetime.
\item Selection changes relative-entropy reconstruction or changes the
source through an actual intervention. A returned geometric record can
then change a permitted preparation; deleting the record after emission
does not undo its stress.
\item Finite-radius corrections have one causal source/geometry account,
with separately controlled readout and collar errors. The small-ball
limit is uniform in the prepared family, rather than obtained by spending
unbounded resources on each ball independently.
\end{enumerate}
\end{conjecture}

The flat Weyl and neutral pulses provide nonzero-frequency source
witnesses for the first order; the scalar two-jet above provides an
explicit propagation and Ward test for the next order. HAFF-D provides
finite source--response--record returns in a supplied JT sector.
Q-RAIF~C provides an interacting quantum source with retained references
and a common classical covariant action~\cite{QRC}. These are different
models and dimensions. None yet proves this conjecture's common
higher-dimensional quantum source or its actual curved modular response.
The missing bridge is substantive: flat CHM covariance does not persist
for arbitrary $g$, quantum reference stress can gravitate, and a finite
positive regulator does not establish a renormalized covariant Ward law.

The bent-slice example is a geometric falsification test. A second is
to insert the flat witness unchanged at order $\theta^4$: unless
$\mathcal D_{h_2}[t_2]=0$ exceptionally, it fails the source Ward identity.
A third compares charged and invariant-algebra relative entropies on
one physical preparation while keeping the complete source fixed.
Any resulting difference is a reconstruction difference unless a new
source operation is demonstrated. These tests preserve the proposed
connection between physical access and geometry without assigning
gravitational stress to information that has merely been forgotten.

\section{Higher multipoles and finite reading}\label{sec:multipoles}

Under a supplied radial Gaussian readout model, the ball coefficients
give conditional information $\log I_k=-4k\log k+O(k)$.
The resolved order consequently grows as $\log N/(4\log\log N)$,
as in Consciousness~III's information-spectrum bound \cite{PaperIII}.
Appendix~\ref{app:prop:multipoles} retains the normalized directions,
noise assumptions, exact Schur complement and proof. This is a
world-source reading model, not a self-description theorem.

\section{Relation to other work, and limits}\label{sec:relation}

\paragraph{Entanglement equilibrium.} Jacobson~\cite{Jacobson2016} compares, in a small
geodesic ball, the variation of the area at fixed volume with the CHM modular energy, and
obtains the Einstein equation from all rest frames at leading order; the coefficient
\(\Omega_{d-1}R^{d+1}/(d(d+2))\) is the \(c_0\) of~\eqref{eq:series}. The present paper adds
the following: the fixed-volume coefficient is the fixed-reference variational multiplier and
follows from full-kernel matching (Proposition~\ref{prop:lambda}); the
comparison is exact at every radius (Theorem~\ref{thm:equivalence}); among functionals of
the intrinsic geometry, internal scales allowed, nonzero-frequency witnesses identify it up to normalization
(Theorem~\ref{thm:uniqueness}); the condition has to be written with modular energy
rather than entropy (Remark~\ref{rem:energy}, Proposition~\ref{prop:entropy}); and the
results say what an observer's partial data fix (Section~\ref{sec:access}).

\paragraph{The first law of causal diamonds.} Jacobson and Visser~\cite{JacobsonVisser2019}
derive, from the covariant phase space of general relativity, a first law relating the
variations of the edge area, the volume of the maximal slice, the cosmological constant and
the matter Hamiltonian of a diamond in a maximally symmetric space.
Theorem~\ref{thm:identity} is its flat-space, linear geometric content, proved without a
field equation.

\paragraph{Holography.} Faulkner et al.~\cite{Faulkner2014} show that the first law for
all boundary balls of a holographic CFT is equivalent to the linearized Einstein equation of
the bulk; HAFF-D uses this for a supplied half-sided modular family~\cite{HAFFD}. There the
geometry is the bulk geometry and the Ryu--Takayanagi formula supplies the geometric
side~\cite{RyuTakayanagi2006}. Here the geometry is that of the spacetime in which the
balls live, no dictionary is used, and \((\mathrm{MC}_\eta)\) takes the place of the
Ryu--Takayanagi formula. The present result is a same-spacetime conditional route, with \((\mathrm{MC}_\eta)\) as its explicit hypothesis.

\paragraph{Non-conformal matter and higher curvature.} For non-conformal fields the
modular Hamiltonian of a ball is not the CHM charge, and Casini, Galante and Myers showed
that for relevant deformations of low dimension the corrections dominate the small-ball
limit~\cite{CGM2016}. The modular results and scalar preparations here use
conformal matter; the leading two-jet geometric test only requires its
stated smooth source. Higher-curvature versions
of entanglement equilibrium use the Wald entropy and a generalized
volume~\cite{BMSV2017}, consistent with Proposition~\ref{prop:failures}(b).

\paragraph{Limits.} The exact finite-radius results are linear in \(h\) about Minkowski space,
with conformal matter and the specified bulk modular datum.
Section~\ref{sec:secondorder} extends the local small-ball law to a geometric
two-jet, with a separate stress consistency input. A physical boundary charge cannot be
discarded when it varies; the scalar example makes both its compensating readout and
the failure of the raw-energy law explicit. The constant \(\eta\) is supplied.
\((\mathrm{MC}_\eta)\) is a hypothesis:
the equivalence theorem identifies its field equation, and nonzero-frequency witnesses
identify its intrinsic response up to normalization. They do not say why the modular energy of a ball should be balanced by
a geometric response at all. Only balls are used.

\section{Access, common geometry and the remaining mechanism}
\label{sec:meaning}
\label{sec:open}

The ball family gives a precise sense in which different perspectives
must agree with one geometry. Given the linear consistency condition,
its intrinsic response is the Einstein tensor. A prepared nonzero-frequency
source identifies that response; homogeneous or zero sources do not.
The metric additionally needs gauge and common gravitational data.
Comparison of rest frames recovers components inaccessible to a single
frame, and finite rapidity or radius sets the precision of that recovery.

Stress-containing selected algebras share this sourced tensor. A coherent
preparation loses relative entropy when restricted to an invariant algebra,
while its stress remains. This distinguishes an observer's reach from a
second physical geometry. The difference matters to Consciousness~III's
finite self-models without implying that a modular energy or a ball algebra
is itself experienced~\cite{PaperIII}.

The two-jet extension makes the next step concrete. An intrinsic ball
alone can confuse a bent slice with a gravitational source; the Gauss
correction removes that error. All small boosted balls then test one
second-order residual, including the quadratic first-order geometry.
The scalar source follows its perturbed propagation and curvature coupling.
It cannot be held flat while its geometry is evolved. These are necessary
parts of a causal selection--source--response mechanism, not its full
derivation.

The central question is why the actual modular data obey the required
geometric consistency at all. HAFF-D addresses physical selection and
returned source response~\cite{HAFFD}; the present paper gives their
higher-dimensional consistency and identification conditions. A common
renormalized quantum source, curved modular response and controlled
finite-radius law remain the high-value bridge in
Conjecture~\ref{conj:secondorder-common}. Curved backgrounds,
non-conformal matter, the origin of $\eta$ and shape-sensitive tidal
readings remain further directions.

\appendix
\section{Linear ball kernels and response identification}
\label{app:linear-details}
\subsection{The small-ball series and exact Fourier transfer}\label{subsec:series}

\begin{proposition}[Small-ball series]\label{prop:series}
Let \(\dT\) be smooth near \(x_0\). For every \(u\) and every \(N\),
\begin{equation}\label{eq:series}
\begin{gathered}
\delta\ev{K_B}=2\pi\sum_{k=0}^{N}c_k\,R^{d+1+2k}\,\Delta_u^k\,\dT_{uu}(x_0)
+O\big(R^{d+3+2N}\big),\\
c_k=\frac{\Omega_{d-1}}{4^k\,k!\,(d/2)_k\,(d+2k)(d+2k+2)},
\end{gathered}
\end{equation}
with \(\Omega_{d-1}=2\pi^{d/2}/\Gamma(d/2)\) and \((a)_k\) the Pochhammer symbol. For a
single Fourier component \(\dT_{\mu\nu}=\mathrm{Re}(\varepsilon_{\mu\nu}e^{ik\cdot x})\),
exactly,
\begin{equation}\label{eq:transfer}
\delta\ev{K_B}=\frac{2\pi V_d}{d+2}\,R^{d+1}\,f\big(R\,|k_u|\big)\,
\mathrm{Re}\big(\varepsilon_{uu}e^{ik\cdot x_0}\big),\qquad
f(z)=\Gamma\!\big(\tfrac d2+2\big)\Big(\frac2z\Big)^{d/2+1}J_{d/2+1}(z),
\end{equation}
where \(|k_u|^2=k\cdot k+(u\cdot k)^2\) is the squared spatial momentum in the frame \(u\),
\(V_d=\Omega_{d-1}/d\) and \(f(z)=1-z^2/(2(d+4))+O(z^4)\).
\end{proposition}

The proof is in Appendix~\ref{app:prop:series}.


The leading coefficient is \(c_0=V_d/(d+2)\). For smooth inputs this is
a finite-order asymptotic expansion; an infinite convergent series requires
additional analyticity. The Fourier formula is exact.
Section~\ref{sec:multipoles} uses these radial coefficients.

\begin{corollary}[Tensor recovery]\label{cor:recovery}
The leading terms of~\eqref{eq:series} for all \(u\) in any nonempty open subset of the unit
hyperboloid determine \(\dT_{\mu\nu}(x_0)\).
\end{corollary}

\begin{proof}
If \(A_{\mu\nu}u^\mu u^\nu=0\) on an open subset of the hyperboloid, it vanishes on all of
it by analyticity, hence on the open cone of timelike vectors by homogeneity, so \(A=0\).
\end{proof}

A single frame gives only \(\dT_{uu}\) at each point; Section~\ref{subsec:singleframe}
shows that this is a real loss, not an artefact of the leading order.


\subsection{The small-ball coefficients}
\label{app:prop:series}
\begin{proof}
The spherical mean of a smooth function over the sphere of radius \(r\) about \(x_0\) in
\(\Sigma_u\) is \(\sum_k r^{2k}\Delta_u^k g(x_0)/(4^kk!(d/2)_k)\) up to \(O(r^{2N+2})\)
(Pizzetti's formula~\cite{Pizzetti1909}), and
\(\int_0^R\frac{R^2-r^2}{2R}r^{d-1+2k}dr=R^{d+1+2k}/((d+2k)(d+2k+2))\). For the Fourier
component, \(\int_{|y|<1}(1-|y|^2)e^{ip\cdot y}d^dy=\pi^{d/2}(2/p)^{d/2+1}J_{d/2+1}(p)\), and
\(e^{ik\cdot x}\) restricted to \(\Sigma_u\) is a plane wave of momentum \(|k_u|\).
\end{proof}
\subsection{Intrinsic distributional classification}
\label{app:lem:tau}
\begin{proof}
Translate the centre to zero. By (F6) write
\(\delta F_B[h]=\sum_{ij}\langle L_{ij},h_{ij}\rangle\), with a symmetric tensor
distribution \(L\) supported in the closed ball. Gauge invariance gives
\(\partial_iL_{ij}=0\), and rotation invariance makes
\(\sigma=\sum_iL_{ii}\) a radial distribution. A constant metric perturbation is
pure gauge, so \(\langle\sigma,1\rangle=0\).

Let \(\Gamma\) be the radial fundamental solution of the spatial Laplacian,
\(\Delta\Gamma=\delta_0\), and set
\(\mu_B=-\Gamma*\sigma/(d-1)\). This distribution is supported in the closed ball.
Indeed, for \(|x|>R\) the function \(y\mapsto\Gamma(x-y)\) is harmonic on a
neighbourhood of the ball. Its spherical average in \(y\) is the constant
\(\Gamma(x)\), by the mean-value property. Pairing with a radial distribution
therefore gives \((\Gamma*\sigma)(x)=\Gamma(x)\langle\sigma,1\rangle=0\).
The argument includes the logarithmic fundamental solution in \(d=2\).

For real \(p\ne0\), orthogonal covariance implies
\(\widehat L_{ij}=a(|p|)\delta_{ij}+b(|p|)p_ip_j\).
Transversality and the trace then give
\[
\widehat L_{ij}(p)
=(|p|^2\delta_{ij}-p_ip_j)
 \frac{\widehat\sigma(p)}{(d-1)|p|^2}
=(|p|^2\delta_{ij}-p_ip_j)\widehat\mu_B(p).
\]
Both sides are Fourier transforms of compactly supported distributions, so equality
extends through zero and identifies the distributions. Thus
\(L_{ij}=(\partial_i\partial_j-\delta_{ij}\Delta)\mu_B\).
Distributional integration by parts proves (d) for arbitrary smooth \(h\).
The compact radial distribution \(\mu_B\) has an even entire Fourier transform
along any coordinate axis, which proves (a). Differentiation at zero proves (b).
For (c), the Gauss equation
\({}^{(d)}\!R+K^2-K_{ij}K^{ij}=2G_{\mu\nu}n^\mu n^\nu\) has \(K_{ij}=O(h)\) about the flat
slice, so at linear order \(\tau_u=2\Glin_{uu}\).
\end{proof}
\subsection{Fixed-reference multiplier and full kernel}
\label{app:prop:lambda}
\begin{proof}
(a) is a one-line differentiation. For (b), the functional depends on \(h\) only through
\(\tau_u\) (Lemma~\ref{lem:tau}), so it can be evaluated on the conformally flat
perturbation \(h_{ij}=2\phi\,\delta_{ij}\), for which \(\tau_u=-2(d-1)\Delta\phi\). With the
spherical means \(\sum_km_ks^{2k}\) of \(\phi\), the linearized geodesic ball gives
\(\delta A=(d-1)\Omega_{d-1}\sum_km_kR^{d-1+2k}\frac{2k}{2k+1}\) and
\(\delta V=\Omega_{d-1}\sum_km_kR^{d+2k}\frac{2k(d-1)}{(2k+1)(d+2k)}\); comparison with
Proposition~\ref{prop:series} gives the ratio.
\end{proof}
\subsection{The nonzero-frequency and Radon identification proof}
\label{app:thm:uniqueness}
\begin{proof}
On any open set of centres, the real and imaginary parts of a nonzero plane wave
are linearly independent. Lemma~\ref{lem:tau}(d) therefore gives
\[
2\eta m_R(q_u)g_{uu}+2\pi\widehat w_R(q_u)t_{uu}=0,\qquad
\widehat w_R(q)=\frac{V_dR^{d+1}}{d+2}f(Rq).
\]
The entire function \(\widehat w_R\) is nonzero. The transfer function \(m_R\)
cannot vanish identically: otherwise \(t_{uu}\) would vanish on an open frame set
and tensor recovery would give \(t=0\).
Exclude the isolated zeros of \(m_R\) and \(\widehat w_R\), and restrict to a
frame patch whose \(q_u\) varies over an interval. There
\[
g_{uu}=a(q_u)t_{uu},\qquad
a(q)=-\frac{\pi\widehat w_R(q)}{\eta m_R(q)}.
\]
Lemma~\ref{lem:quadric} gives \(g=a_0t\); the scalar is nonzero on this patch.
The entire identity
\(2\eta a_0m_R+2\pi\widehat w_R=0\) follows by analytic continuation.
Set \(c=2\pi/(\eta a_0)\) to obtain the normalization.
The real, reflection-invariant functional has a real transfer function on real momenta,
so \(a_0\) and \(c\) are real.
A compactly supported tensor distribution is determined by its Fourier transform,
so equality on plane waves identifies the functional on all smooth inputs.
The fixed nonzero amplitudes fix \(a_0\), hence the same \(c\), at every other
tested radius.

For the measure extension, dominate all tensor components by the locally finite measure
\(\mu=\sum_{a,b}(|\widehat{\dT}_{ab}|+|\widehat G_{ab}|)\), and write their
Radon--Nikodym densities as \(t_{ab}(k),g_{ab}(k)\).
Lemma~\ref{lem:tau}(d) represents the centre dependence as convolution with a compact
distribution on the slice. Its Fourier multiplier is \(m_R(q_u(k))\), a smooth
function of \(k\) whose real-axis growth is polynomial. Multiplication therefore
preserves the tempered Radon class. Conservation holds for these
densities almost everywhere. Fourier-transform consistency for a countable dense set
of frames in the open frame set. Intersect their full-measure sets: at each remaining
\(k\), the displayed coefficient equation holds on the dense frame set, hence on
the entire open set by continuity in \(u\). The source hypothesis supplies one
such \(k\ne0\) with \(t(k)\ne0\). The amplitude proof at that \(k\) identifies
\(m_R=-c\widehat w_R/2\).

Substitute this identity back into the measure equation. For almost every \(k\),
\(\widehat w_R(q_u)[g(k)-(2\pi/(c\eta))t(k)]_{uu}=0\) on the open frame set.
If \(k\ne0\), \(q_u\) varies over an interval and the entire function
\(\widehat w_R\) has only isolated zeros. The tensor contraction therefore vanishes
on an open frame patch and tensor recovery sets the bracket to zero.
At \(k=0\), \(\widehat w_R(0)>0\) gives the same conclusion directly.
Thus the tensor measures, and hence the distributions, obey the field equation.
This uses densities with respect to a common measure, without dividing distributions.
At any additional tested radius repeat the argument with that radius's own exceptional
null set. The already established field equation fixes the same constant \(c\), since
the source is nonzero. No common full-measure set for uncountably many radii is required.
\end{proof}
\subsection{The zero-frequency converse}
\label{app:cor:spectraliff}
\begin{proof}
Sufficiency is Theorem~\ref{thm:uniqueness}. A Radon measure supported at zero is a
multiple of \(\delta_0\), so the remaining nonzero source is constant. In the
common-measure proof above, its atom forces \(m_R(0)\ne0\). At almost every
\(k\ne0\) the source density is zero and consistency gives
\(m_R(q_u)g_{uu}=0\). Since \(m_R\) is nonzero and entire, varying the frame
forces \(g(k)=0\). The response tensor is therefore constant as well. Adding
\(\ell^2\int_Bw_R\Delta_u\tau_u\), with \(\ell\ne0\), leaves these data unchanged
but changes the transfer function by \(-\ell^2q^2\widehat w_R(q)\), which is
independent of the original transfer function because it vanishes at zero.
\end{proof}
\section{Prepared sources, bounded pulses and finite readings}
\label{app:source-details}
\subsection{The coherent source and its Radon measure}
\label{app:prop:prepared}
\begin{proof}
The Weyl commutator shifts the field by the classical solution \(\theta\varphi_f\);
normal ordering therefore gives~\eqref{eq:preparedsource} and
\(\delta\ev{\sum_a:\!\phi_a^2\!:}=\theta^2\varphi_f^2\).
Since \(\Box\varphi_f=0\), differentiation gives conservation and zero trace.
Finite propagation makes its Cauchy data compactly supported on every fixed time slice.
The improvement in \(T_{00}\) is a spatial divergence, so
\[
\int_{\mathbb R^d}\dT_{00}\,d^dx
=\frac{\theta^2}{2}\int_{\mathbb R^d}
\big[(\partial_t\varphi_f)^2+|\nabla\varphi_f|^2\big]\,d^dx>0.
\]
Zero energy would make \(\varphi_f\) constant and, by compact spatial support, zero,
contrary to the nonzero one-particle class.

Up to Fourier convention, the spacetime transform of \(\varphi_f\) is
\(\operatorname{sgn}(k^0)\delta(k^2)\widehat f(k)\,d^{d+1}k\).
Its two mass-shell sheets carry densities proportional to
\(\widehat f(\pm|p|,p)\,d^dp/(2|p|)\).
Near \(p=0\) the radial total-variation integral is \(O(r^{d-2}dr)\), integrable for
\(d\ge2\); at infinity the restricted Schwartz function decreases faster than any power.
Thus this is a finite complex measure with every polynomial moment finite.
Derivatives multiply it by momenta. Products of these smooth fields have Fourier
measures given by convolution; the improvement differentiates the convolution twice.
These operations retain finite variation and every polynomial moment, proving the
assertion for \(\widehat{\dT}\).
If that measure were supported only at zero, \(\dT\) would be a constant tensor.
Compact spatial support would force it to vanish, contradicting its positive total energy.
Finally, a point on the edge of \(B\) is spacelike to every point in a compact subset
of \(D(B)\), so \(\varphi_f\) vanishes there and~\eqref{eq:preparededge} is zero.
\end{proof}
\subsection{The invariant quadratic source}
\label{app:prop:neutral}
\begin{proof}
The CCR and \([X_a,Q_f]=0\) give the exact transformation
\(U_\theta^*\phi_a(x)U_\theta=\phi_a(x)+\theta v_f(x)X_a\).
Applying it to the vacuum two-point function and the point-split stress tensor gives
\eqref{eq:neutralsource} and the boundary formula. Internal rotations preserve
\(X_1^2+X_2^2\); injectivity of \(\arctan\) on the spectrum implies that its bounded
functional calculus contains \(U_\theta\). The stable-seed theorem then proves the
hull assertions.

Both \(v_f\) and \(H_f\) are smooth real massless solutions. Their Fourier measures
have the same rapidly decreasing mass-shell densities as in
Proposition~\ref{prop:prepared}, with different phases on the two sheets. For
\(d\ge2\) they have finite variation and every polynomial moment. The products and
derivatives in~\eqref{eq:neutralsource} preserve these properties. Every term contains
\(v_f\) or a derivative of it, so finite propagation gives compact spatial support.
Polarization of the free-field identities proves conservation and zero trace.

The complex solution \(\omega(X_a\phi_a(x))\) has only one frequency sign.
Its real and imaginary parts have zero cross term in the integrated classical energy: in spatial Fourier variables
the two amplitudes differ by multiplication by \(i\). Since \(v_f\) is a real
multiple of the imaginary part and \(H_f\) is its real part,
\(\int t_{00}[v_f,H_f]\,d^dx=0\). The remaining energy is the displayed positive
quadratic term. A Fourier measure supported only at zero would make the source
constant, which contradicts compact spatial support and positive total energy.
\end{proof}
\subsection{The bounded accessible pulse}
\label{app:prop:bounded-source}
\begin{proof}
The weak integral belongs to \(\mathcal N\). Choose normalized nonnegative
smearings whose supports shrink to a point \(t_0\in(0,\tau)\). Then
\(B\Omega\to\alpha_{t_0}(C)\Omega=e^{it_0H}C\Omega\).
Separation keeps \(C\Omega\), hence this limit, off \(\mathbb C\Omega\), so
sufficiently narrow smearings remain non-scalar.
Translation smoothing gives a bounded commutator
\([H,B]=i\int f'(t)\alpha_t(C)dt\).
Consequently \(U_\theta\) preserves \(\operatorname{Dom}H\), with
\(\|[H,U_\theta]\|\le |\theta|\|[H,B]\|\), and its vacuum excitation has
finite energy. For \(0<|\theta|\|B\|<\pi\), the exponential is injective on the
bounded spectrum of \(B\), hence \(U_\theta\) is non-scalar. If its energy
were zero, the unique ground vector would give \(U_\theta\Omega=z\Omega\),
and separation would force \(U_\theta=zI\), a contradiction.
The upper bound follows from
\(E_\theta=\langle U_\theta\Omega,[H,U_\theta]\Omega\rangle\).
The bounded commutator makes the exponential expansion valid in the graph norm
of \(H\). Since \(H\Omega=0\), the double commutator gives the quadratic term;
the spectral theorem gives its integral and strict positivity. Put
\(R_\theta=(U_\theta-I-i\theta B)\Omega\).
The exponential integral formula gives
\(\|R_\theta\|\le\theta^2\|B\|^2/2\) and
\(\|HR_\theta\|\le\theta^2\|B\|\|K\|\).
Expanding the energy of \(i\theta B\Omega+R_\theta\) gives the stated remainder.
\end{proof}
\subsection{Strict entropy loss and the regulated coefficient}
\label{app:prop:entropy}
\begin{proof}
(a) The equality is~\eqref{eq:relent-unitary}; it is finite because a smooth coherent
excitation has finite modular energy. Here the compactly supported preparation is strictly
local to \(D(B)\), so its surface variation~\eqref{eq:preparededge} vanishes for this ball.
Take a real test function \(g\) in \(D(B)\) with
nonzero commutator function \(\Delta(f,g)\), and \(x=\arctan\phi_1(g)\in\mathcal A(D(B))\). The
rotation by \(\pi\) maps \(\phi_1\) to \(-\phi_1\) and commutes with \(E_B\), and \(\arctan\)
is odd, so \(E_B(x)=0\). In \(\omega_W\) the field \(\phi_1(g)\) has the vacuum Gaussian law
shifted by \(\theta\Delta(f,g)\neq0\), so \(\omega_W(x)\neq0=\omega_W(E_B(x))\). The state does
not factor through \(E_B\), and by Petz's criterion for sufficient
subalgebras~\cite{Petz1986} the inequality is strict. (b) The coherent stress expectation is quadratic in its classical displacement.
Rotation by \(\pi\) makes the restricted state even in \(\theta\).
For a finite faithful density matrix, differentiating the logarithm gives
\[
g_{\mathrm{BKM},\rho_0}(\chi,\chi)=
\sum_{a,b}|\chi_{ab}|^2
 \frac{\log p_a-\log p_b}{p_a-p_b},
\]
with the diagonal value \(1/p_a\); the first entropy derivative vanishes because
\(\operatorname{Tr}\chi=0\). This proves the regulated assertion.
\end{proof}
\subsection{Radial multipoles and finite reading}

The higher terms of~\eqref{eq:series} are the radial multipoles of the local source, read
through the modular weight. Consciousness~III states a resolution
order for any family of directions whose information falls as \(e^{-\alpha j\log j}\), and
notes that a physical instance needs normalized directions, a fluctuation kernel and a
permitted measurement~\cite{PaperIII}. The ball family supplies the directions and the
signal; the noise remains to be supplied.

\begin{proposition}[Fisher spectrum of the multipoles]\label{prop:multipoles}
Write the normalized reading of a ball as
\begin{gather*}
y(R)=\frac{\delta\ev{K_B}}{2\pi V_dR^{d+1}/(d+2)}=\sum_{k\ge0}a_k\Big(\frac RL\Big)^{2k}\theta_k,\\
\theta_k=L^{2k}\Delta_u^k\dT_{uu}(x_0),\qquad
a_k=\frac{d(d+2)}{4^kk!\,(d/2)_k\,(d+2k)(d+2k+2)},
\end{gather*}
exact when the analytic ball-average series converges on the radius range used.
Suppose a record consists of \(n\) readings at radii drawn
independently and uniformly from \((0,\rho L]\), each with an independent Gaussian error of
variance \(\sigma^2\). The Fisher information per record about \(\theta_k\), with
\(\theta_0,\dots,\theta_{k-1}\) unknown and all higher coefficients fixed and known
(or with the series truncated at order \(k\)), is
\[
I_k=\frac{n}{\sigma^2}\,a_k^2\rho^{4k}\,\frac1{4k+1}\prod_{j<k}\frac{[2(k-j)]^2}{[2(k+j)+1]^2},
\qquad \log I_k=-4k\log k+O(k),
\]
and \(\rho\) enters only at order \(k\).
\end{proposition}

\subsection{The radial Fisher Schur complement}
\label{app:prop:multipoles}
\begin{proof}
The Fisher matrix is \(F_{jk}=(n/\sigma^2)a_ja_k\rho^{2j+2k}/(2j+2k+1)\), a diagonally scaled
Cauchy matrix with entries \(1/(x_j+x_k)\), \(x_j=2j+\frac12\). The information about
\(\theta_k\) with the lower parameters unknown is the last Schur complement of the leading
\((k+1)\times(k+1)\) block, for which the Cauchy determinant gives the product. Stirling's
formula gives \(\log a_k=-2k\log k+O(k)\), and the product is \(e^{O(k)}\).
\end{proof}

\subsection{Numerical checks and applications}
The closed form agrees with a Cholesky factorization at eighty-digit precision for
\(k<14\) and \(d=2,3,4\). A readout weight comparable to a power of \(R\), in
the radius distribution or in the noise variance, changes the Schur complements by factors
\(e^{O(k)}\) and leaves the leading term. With Consciousness~III's tolerance \(c\), the
number of multipoles resolved from \(N\) records is
\(K(N)\sim\log N/(4\log\log N)\); for \(d=3\), \(\rho=1\) and \(c=1\) one finds \(K=8\) at
\(\log N=100\) and \(K=55\) at \(\log N=1000\). The factorial suppression is a property of
the modular weight, so the resolution exponent \(\alpha=4\) has a physical signal behind it.
The readout model, including its noise, is supplied, and the directions belong to the world
as read by the observer, not to the observer's model of itself.

\section{Regularity of the quartic entropy law}\label{app:bkm}

For a displaced charged thermal mode, put $z=e^{-b}$,
$p_n=(1-z)z^n$. Phase averaging gives the tangent
$\chi_n=|\alpha|^2p_n[(1-z)^2n/z-(1-z)]$.
The geometric variance $z/(1-z)^2$ yields
\begin{equation}\label{eq:modebkm}
 g_{\mathrm{BKM}}(\chi,\chi)
 =\sum_n\chi_n^2/p_n
 =|\alpha|^4(1-z)^2/z
 =4|\alpha|^4\sinh^2(b/2).
\end{equation}
Thus finite modular energy $\sum_jb_j|\alpha_j|^2$ does not alone
control the diagonal BKM coefficient.

For finitely many independent modes of one charge, the invariant
coherences also contribute. With $m=n+e_j-e_l$, expansion of the
displacement gives
\[
 \chi_{mn}=\alpha_j\overline\alpha_l
 \sqrt{(n_j+1)n_l}\,p_n(1-z_j)(1-z_l)/z_l.
\]
Using $p_m/p_n=z_j/z_l$ in the spectral BKM kernel and averaging
$(n_j+1)n_l$ gives
\begin{equation}\label{eq:multimodebkm}
 g_{\mathrm{BKM}}(\chi,\chi)
 =\sum_{j,l}|\alpha_j|^2|\alpha_l|^2
 (1-z_j)(1-z_l)\frac{\log z_j-\log z_l}{z_j-z_l},
\end{equation}
with coincident value $(1-z_j)^2/z_j$.
Opposite charges add the ordered neutral-pair sum with kernel
$(1-z_j)(1-z_l)(b_j+b_l)/(1-z_jz_l)$.
The mean-value theorem bounds the same-charge kernel by
$e^{b_j}+e^{b_l}$; the neutral-pair kernel is at most $b_j+b_l$.
Consequently, writing
$A=\sum_j|\alpha_j|^2$, $E=\sum_je^{b_j}|\alpha_j|^2$ and
$L=\sum_jb_j|\alpha_j|^2$, the full regulated tangent obeys
$g_{\mathrm{BKM}}\le2A(E+L)$.
The full mode expansion is retained as
\nolinkurl{CII_BKM_mode_calculation.tex} in the supplementary source record.

These estimates control the tangent, not the entropy remainder.
For the continuum ball net, a quartic law still requires its modular
BKM norm and uniform remainder control as occupation, mode and volume
regulators are removed. Compact smooth spacetime smearing has not
been shown here to supply the exponential moment $E$.
The exact strict relative-entropy inequality is independent of this
open regularity question.

\section{Explicit two-jet geometry and source bookkeeping}
\label{app:twojet}

All index moves in the following coefficients use $\eta$, and
antisymmetrization is written out rather than hidden in a factor.
For $g=\eta+\epsilon h+\epsilon^2j$, the connection coefficients are
\begin{align*}
 \Gamma^{(1)\lambda}_{\mu\nu}[h]
 &=\tfrac12\eta^{\lambda\rho}
  (\partial_\mu h_{\nu\rho}+\partial_\nu h_{\mu\rho}
                               -\partial_\rho h_{\mu\nu}),\\
 \Gamma^{(2)\lambda}_{\mu\nu}[j,h]
 &=\Gamma^{(1)\lambda}_{\mu\nu}[j]
   -h^\lambda{}_{\rho}\Gamma^{(1)\rho}_{\mu\nu}[h].
\end{align*}
Define the quadratic part with $j=0$ by $\gamma_2=\Gamma^{(2)}[0,h]$.
Then
\begin{align*}
 r^{(2)}_{\mu\nu}[h,h]
 ={}&\partial_\lambda\gamma_2^\lambda{}_{\mu\nu}
       -\partial_\nu\gamma_2^\lambda{}_{\mu\lambda}
       +\Gamma^{(1)\lambda}_{\lambda\rho}\Gamma^{(1)\rho}_{\mu\nu}
       -\Gamma^{(1)\lambda}_{\nu\rho}\Gamma^{(1)\rho}_{\mu\lambda},\\
 G^{(2)}_{\mu\nu}[h,h]
 ={}&r^{(2)}_{\mu\nu}
       -\tfrac12\eta_{\mu\nu}
          (\eta^{\alpha\beta}r^{(2)}_{\alpha\beta}
                       -h^{\alpha\beta}R^{\mathrm{lin}}_{\alpha\beta}[h])
       -\tfrac12h_{\mu\nu}R^{\mathrm{lin}}[h].
\end{align*}
These coefficients include the inverse-metric correction. The full
quadratic Einstein coefficient is $\Glin[j]+G^{(2)}[h,h]$, not
$\Glin[j]$ alone. For a first-order vacuum wave, the second-order
vacuum equation is $\Glin[j]=-G^{(2)}[h,h]$; interpreting its right
side as a gravitational energy tensor requires a gauge/averaging convention.

For the scalar define $s=v^2$ and
\[
 C_h^\rho=\eta^{\alpha\beta}\Gamma^{(1)\rho}_{\alpha\beta},\qquad
 \delta\Box_h s=-h^{\alpha\beta}\partial_\alpha\partial_\beta s
                                      -C_h^\rho\partial_\rho s.
\]
Varying the first three terms of Eq.~\eqref{eq:curvedscalar} at fixed $v$
gives the fully explicit quantity
\begin{align*}
 \delta_g t_{h,\mu\nu}[v]={}&-\tfrac12h_{\mu\nu}(\partial v)^2
       +\tfrac12\eta_{\mu\nu}h^{\alpha\beta}
                                 \partial_\alpha v\partial_\beta v\\
 &+\xi\bigl(h_{\mu\nu}\Box s+\eta_{\mu\nu}\delta\Box_hs
                      +\Gamma^{(1)\rho}_{\mu\nu}\partial_\rho s\bigr).
\end{align*}
At a normal-coordinate point this vanishes when $h$ and its first
derivatives vanish. The remaining $\xi\Glin[h]v^2$ does not, which
is the nonzero-source calculation in Section~\ref{sec:secondorder}.
The identity
$\nabla^\mu[(g_{\mu\nu}\Box-\nabla_\mu\nabla_\nu)s]
=-R_{\nu\rho}\nabla^\rho s$, together with the Bianchi identity,
reduces the full scalar stress divergence to its scalar equation times
$\partial_\nu\Phi$. Its two-jet expansion gives
Eq.~\eqref{eq:secondward}, including the inverse metric and both connection
terms. Counterterms cancel between two states only with one common
metric and renormalization convention; the reference source remains
in a total gravitational equation.

A finite algebraic check uses
$g=e^{2\epsilon f}\eta$,
$\Phi=\sqrt\epsilon e^{-(d-1)\epsilon f/2}v$,
$f=t^2+x^2/3+tx/5$ and $v=1+2t$.
Direct rational-polynomial expansion through $\epsilon^2$ in spacetime
dimensions $3,4,5$ gives the conformal wave equation, vanishing trace
and covariant divergence. Dropping $\xi G\Phi^2$ fails the divergence
test. This tests the source identity on a prescribed conformal metric,
not a solution of the coupled Einstein equations, and does not supply
quantum renormalization or an actual curved modular Hamiltonian.

\newpage
\addcontentsline{toc}{section}{References}
\begin{thebibliography}{99}

\bibitem{PaperI}
S.~Liu, \emph{Geometry Relative to a Cut: Selection, Generation and Modular Inseparability} (Consciousness~I), (2026), series record DOI:
10.5281/zenodo.19607387.

\bibitem{PaperIII}
S.~Liu, \emph{Finite Self-Models and the Explanatory Gap: What a Continuing Observer Can
Know About Itself} (Consciousness~III), (2026), series record DOI:
10.5281/zenodo.19607724.

\bibitem{HAFFD}
S.~Liu, \emph{Selection, Maintenance and Gravitational Response:
Modular Constraints and a Physical Return}, (2026), series
record DOI: 10.5281/zenodo.18388881.

\bibitem{FoundationI}
S.~Liu, \emph{Foundation I: From Source to Manifestation. Accessible Algebras from a
Source, Its Noise and a Readable Seed}, (2026), series record DOI:
10.5281/zenodo.19627786.

\bibitem{QRA}
S.~Liu, \emph{Lorentzian Read-out from Local Unitary Propagation: Clifford Relations, Causal
Cones and Spectral Mass} (Q-RAIF~A), (2026), series record DOI:
10.5281/zenodo.18525876.

\bibitem{QRC}
S.~Liu, \emph{Retained Quantum History, Physical Sources and Receiver Renewal in Q-RAIF} (Q-RAIF~C), (2026), series record DOI:
10.5281/zenodo.18528934.

\bibitem{CasiniHuertaMyers2011}
H.~Casini, M.~Huerta and R.~C.~Myers, \emph{Towards a derivation of holographic
entanglement entropy}, JHEP \textbf{1105} (2011) 036, arXiv:1102.0440.

\bibitem{CasiniMazzitelliTeste2015}
H.~Casini, F.~D.~Mazzitelli and E.~Test\'e, \emph{Area terms in entanglement entropy},
Phys.\ Rev.\ D \textbf{91} (2015) 104035, arXiv:1412.6522.

\bibitem{Herzog2014}
C.~P.~Herzog, \emph{Universal thermal corrections to entanglement entropy for
conformal field theories on spheres}, JHEP \textbf{1410} (2014) 028,
arXiv:1407.1358.

\bibitem{LLPS2015}
J.~Lee, A.~Lewkowycz, E.~Perlmutter and B.~R.~Safdi,
\emph{R\'enyi entropy, stationarity, and entanglement of the conformal scalar},
JHEP \textbf{1503} (2015) 075, arXiv:1407.7816.

\bibitem{HerzogNishioka2016}
C.~P.~Herzog and T.~Nishioka, \emph{The edge of entanglement: getting the boundary
right for non-minimally coupled scalar fields}, JHEP \textbf{1612} (2016) 138,
arXiv:1610.02261.

\bibitem{DonnellyWall2012}
W.~Donnelly and A.~C.~Wall, \emph{Do gauge fields really contribute negatively to
black hole entropy?}, Phys.\ Rev.\ D \textbf{86} (2012) 064042, arXiv:1206.5831.

\bibitem{HislopLongo1982}
P.~D.~Hislop and R.~Longo, \emph{Modular structure of the local algebras associated with the
free massless scalar field theory}, Commun.\ Math.\ Phys.\ \textbf{84} (1982) 71--85.

\bibitem{BGL1993}
R.~Brunetti, D.~Guido and R.~Longo, \emph{Modular structure and duality in conformal quantum
field theory}, Commun.\ Math.\ Phys.\ \textbf{156} (1993) 201--219.

\bibitem{Araki1976}
H.~Araki, \emph{Relative entropy of states of von Neumann algebras}, Publ.\ RIMS
\textbf{11} (1976) 809--833.

\bibitem{Wiesbrock1993}
H.-W.~Wiesbrock, \emph{Half-sided modular inclusions of von Neumann algebras}, Commun.\
Math.\ Phys.\ \textbf{157} (1993) 83--92.

\bibitem{Borchers1992}
H.-J.~Borchers, \emph{The CPT-theorem in two-dimensional theories of local observables},
Commun.\ Math.\ Phys.\ \textbf{143} (1992) 315--332.

\bibitem{Pizzetti1909}
P.~Pizzetti, \emph{Sulla media dei valori che una funzione dei punti dello spazio assume
alla superficie di una sfera}, Rend.\ Lincei \textbf{18} (1909) 182--185.

\bibitem{Jacobson2016}
T.~Jacobson, \emph{Entanglement equilibrium and the Einstein equation}, Phys.\ Rev.\ Lett.\
\textbf{116} (2016) 201101, arXiv:1505.04753.

\bibitem{IyerWald1994}
V.~Iyer and R.~M.~Wald, \emph{Some properties of Noether charge and a proposal for
dynamical black hole entropy}, Phys.\ Rev.\ D \textbf{50} (1994) 846--864,
arXiv:gr-qc/9403028.

\bibitem{JacobsonVisser2019}
T.~Jacobson and M.~R.~Visser, \emph{Gravitational thermodynamics of causal diamonds in
(A)dS}, SciPost Phys.\ \textbf{7} (2019) 079, arXiv:1812.01596.

\bibitem{Faulkner2014}
T.~Faulkner, M.~Guica, T.~Hartman, R.~C.~Myers and M.~Van~Raamsdonk, \emph{Gravitation from
entanglement in holographic CFTs}, JHEP \textbf{1403} (2014) 051, arXiv:1312.7856.

\bibitem{RyuTakayanagi2006}
S.~Ryu and T.~Takayanagi, \emph{Holographic derivation of entanglement entropy from
AdS/CFT}, Phys.\ Rev.\ Lett.\ \textbf{96} (2006) 181602.

\bibitem{CGM2016}
H.~Casini, D.~A.~Galante and R.~C.~Myers, \emph{Comments on Jacobson's ``Entanglement
equilibrium and the Einstein equation''}, JHEP \textbf{1603} (2016) 194, arXiv:1601.00528.

\bibitem{BMSV2017}
P.~Bueno, V.~S.~Min, A.~J.~Speranza and M.~R.~Visser, \emph{Entanglement equilibrium for
higher order gravity}, Phys.\ Rev.\ D \textbf{95} (2017) 046003, arXiv:1612.04374.

\bibitem{Takesaki1972}
M.~Takesaki, \emph{Conditional expectations in von Neumann algebras}, J.\ Funct.\ Anal.\
\textbf{9} (1972) 306--321.

\bibitem{Petz1986}
D.~Petz, \emph{Sufficient subalgebras and the relative entropy of states of a von Neumann
algebra}, Commun.\ Math.\ Phys.\ \textbf{105} (1986) 123--131.

\bibitem{LashkariVanRaamsdonk2016}
N.~Lashkari and M.~Van~Raamsdonk, \emph{Canonical energy is quantum Fisher information},
JHEP \textbf{1604} (2016) 153, arXiv:1508.00897.

\bibitem{FHHPRV2017}
T.~Faulkner, F.~M.~Haehl, E.~Hijano, O.~Parrikar, C.~Rabideau and M.~Van~Raamsdonk,
\emph{Nonlinear gravity from entanglement in conformal field theories}, JHEP \textbf{1708}
(2017) 057, arXiv:1705.03026.

\end{thebibliography}

\end{document}
